\documentclass[10pt]{extarticle}

% ---------- Packages ----------
\usepackage{fancyhdr}
\usepackage{geometry}
\usepackage{hyperref}
\usepackage{xcolor}
\usepackage{amsmath, amsfonts, bm}
\usepackage{subcaption}
\usepackage{graphicx}
\usepackage{tikz}
\usepackage{nomencl}
\usepackage{array}
\usepackage{pdflscape}
\usepackage[export]{adjustbox}
\usepackage[utf8]{inputenc}
\usepackage{parskip}

\usetikzlibrary{shapes.geometric}
\usetikzlibrary{arrows.meta, arrows}

% ---------- Page setup ----------
\geometry{
    a4paper,
    total={170mm,257mm},
    left=20mm,
    top=20mm
}

% ---------- Headers and footers ----------
\pagestyle{fancy}
\pagenumbering{arabic}

\fancyhead[L]{Rocket Physics by Yaroslav Knyazkov}
\fancyfoot[C]{}
\fancyfoot[R]{\thepage}

% ---------- Hyperlinks ----------
\hypersetup{
    colorlinks=true,
    linkcolor={blue!50!black},
    citecolor={blue!50!black},
    urlcolor={blue!80!black}
}

% ---------- Equation numbering ----------
\numberwithin{equation}{subsection}

% ---------- Custom commands ----------
\newcommand{\todo}[1]{%
    \textcolor{red}{(#1)}%
}
% AI-written text is blue; the author's own text stays black. The author reviews the document
% part by part and turns a reviewed part black by removing its wrapper.
%   \ai{...}               inline or multi-paragraph (\newcommand is \long by default)
%   \begin{aitext}...\end{aitext}   whole sections / \input files
% \todo stays red inside either: its own \textcolor overrides the surrounding blue.
\newcommand{\ai}[1]{{\color{blue}#1}}
\newenvironment{aitext}{\color{blue}}{}

\begin{document}

\tableofcontents

\newpage
\begin{aitext}
\section{Current assumptions}\label{sec:assumptions}
This section collects, in one place, every simplification the current derivation and the
simulator rely on. Each is expected to be relaxed in a later phase; the relevant model or backlog
item is named where one exists, so a reader can tell ``current shortcut'' from ``permanent
choice'' at a glance.
\begin{itemize}
    \item \textbf{Planar (2D) flight, no crosswind.} Motion is confined to the vertical plane
    fixed by the launch azimuth $\chi$; the out-of-plane wind component is computed and discarded
    (\autoref{sec:wind}). No vertical wind either: wind velocity is horizontal only,
    updrafts/downdrafts are not modelled and are deliberately left for later, pending research
    into which model is actually worth the added data dependency (\autoref{sec:backlog}).
    \item \textbf{Gravity direction, parallel by default, with two non-parallel models
    available.} $\hat{z}_g$ can rotate with downrange distance, either treating the Earth as a
    sphere (radial) or as the WGS84 ellipsoid via its osculating sphere in the flight azimuth
    (elliptical) --- both in \autoref{sec:gravity} > Direction. Parallel, $\hat{z}_g\equiv
    \hat{z}_0$, is still the default, so every result produced with it is reproduced unchanged.
    Once a non-parallel model is enabled, the ground itself is curved
    (\autoref{eq:surface_point}) and the true altitude/impact condition follow
    \autoref{eq:h_true} rather than $h=-z$, $z=0$.
    \item \textbf{Earth rotation carried only in part, and in the fixed ground frame.} The
    \emph{static} effect is present --- the \emph{magnitude} of $g$ includes the centrifugal
    contribution through the WGS84 model (\autoref{sec:gravity}) --- and so is the in-plane
    Coriolis term in the translational equations of motion (\autoref{sec:noninertial}). What is
    dropped is the out-of-plane Coriolis component, which this 2D model has no degree of freedom
    to receive, and the residual centrifugal term, derived there but left unimplemented as some
    three orders of magnitude smaller. Separately, the decomposition of $\vec{\Omega}$ is frozen
    at the launch site's latitude and heading in $\mathcal{F}_0$ and does \emph{not} follow a
    rotating $\hat{z}_g$; the two corrections stay deliberately independent
    (\autoref{sec:noninertial}, closing caveat).
    \item \textbf{Constant moment arms (not constant mass).} Mass, centre of gravity and moment of
    inertia are \emph{not} held fixed: $m(t)$, $l_g(t)$ and $I(t)$ come from the motor model
    of \autoref{sec:mass} and are carried by \autoref{sec:eom}. What is still held fixed is one
    step further out --- how the other two application points travel relative to a moving centre
    of gravity. The CoG-to-$D$ arm $l_1=l_D-l_B$ never responds to CoG travel, and how far it
    responds to the flight state at all depends on the selected aerodynamic model
    (\autoref{sec:aero}): under the table model it varies with $(\alpha,M)$, while the constant and
    linearized models hold it fixed outright. The CoG-to-$C$ thrust arm is a constant under all
    three.
    \item \textbf{Thrust axis: fixed cant, no commanded control.} The nozzle cant $\delta$
    (\autoref{eq:thrust_frame}) and the lateral offsets $u_B,u_C$ (\autoref{sec:geometry}) are all
    carried as constant inputs --- as-built manufacturing quantities to be swept, not functions of
    the flight state. $\delta=u_B=u_C=0$ is the default and reproduces every earlier result
    unchanged. A \emph{commanded} deflection --- a gimbal and the control law driving it, the
    only pitch authority a rocket has off the rail --- is deliberately not modelled yet
    (\autoref{sec:backlog}).
    \item \textbf{Launch rail, recovery and air brake: modelled, still simplified.} A launch rail
    of length $L_{rail}$ constrains translation to one fixed direction and freezes $\theta$ until
    the rocket clears it, with no friction and no tip-off (\autoref{sec:rail}); a recovery canopy,
    once deployed at apogee or on a timer, adds a drag force opposing the inertial relative-wind
    vector and likewise freezes $\theta$, modelling neither its own opening transient nor the
    body's pendulum swing underneath it (\autoref{sec:recovery}); and an air brake modulates axial
    drag during the unpowered coast, as a switchable extra drag area standing alone rather than a
    surface derived from an airframe --- the wing-body treatment of \autoref{sec:backlog} ---
    affecting $A_x$ only, neither the normal force nor the centre of pressure
    (\autoref{sec:airbrake}). Of the two things simplified in that control loop, only one still is.
    The \emph{law} is continuous and memoryless, re-evaluated afresh from the current flight state
    at every solver call, with no integral or derivative term and no notion of a previous sample
    (\autoref{sec:airbrake_law}). The \emph{integration} no longer rules out a law with memory: an
    opt-in fixed-rate, zero-order-hold mode (\autoref{sec:sampled}) evaluates a sampled quantity
    once per tick and holds it until the next, so that a value surviving from one tick to the next
    is well defined. A controller with memory is therefore expressible but unwritten --- what is
    missing is the law, not somewhere to run it (\autoref{sec:backlog}).
    \item \textbf{Turbulence: a real spectrum, driven at a fixed reference airspeed.} The wind
    fluctuation about the mean profile is the MIL-spec Dryden spectrum, realised as a frozen
    random-phase sum (\autoref{sec:wind} > Turbulence) --- with turbulence off by default, so every
    result produced without it is reproduced unchanged. Two simplifications ride with it, both
    deliberate: the spectrum is driven by a \emph{fixed} reference airspeed rather than the
    rocket's own, which fixes where the spectrum's knee sits in frequency while leaving the
    delivered intensity exact (\autoref{sec:backlog}); and there is no vertical component, for the
    same reason the mean profile has none.
\end{itemize}
\end{aitext}

\newpage
\begin{aitext}
\section{Rocket 2D}\label{sec:rocket2d}
\autoref{sec:notation} lists every symbol used in this report together with the identifier the
simulator uses for it.

\subsection{Frames}\label{sec:frames}
\begin{itemize}
    \item $\mathcal{F}_0$ --- ground/launch frame, origin at the launch point $O^*$.
    \item $\mathcal{F}_g$ --- ground/gravity frame: $\hat{x}_g$ downrange, $\hat{z}_g$ down, along
    gravity.
    \item $\mathcal{F}_r$ --- rocket (body) frame: $\hat{x}_r$ along the nose axis, forward.
    \item $\mathcal{F}_t$ --- thrust frame: $\hat{x}_t$ along the line of thrust.
\end{itemize}
$\mathcal{F}_0$ is fixed at the launch point $O^*$, with $\hat{x}_0$ horizontal (downrange) and
$\hat{z}_0$ down. The rocket's state is always integrated in this frame, regardless of which gravity
or ground model is switched on --- it is the one fixed thing in the whole document. The
translational state is the position $\vec{x}_r=x\,\hat{x}_0+z\,\hat{z}_0$, the velocity
$\vec{v}_r=\mathrm{d}\vec{x}_r/\mathrm{d}t$ and the acceleration
$\vec{a}_r=\mathrm{d}\vec{v}_r/\mathrm{d}t$, all resolved in $\mathcal{F}_0$.

$\mathcal{F}_g$ carries the direction gravity actually points, $\hat{z}_g(\varphi)$, a function of
the downrange arc angle $\varphi$ introduced in \autoref{sec:gravity}. It coincides with
$\mathcal{F}_0$ by default, $\hat{x}_g(0)\equiv\hat{x}_0$ and $\hat{z}_g(0)\equiv\hat{z}_0$, and
only rotates away from it once a non-parallel direction model is enabled --- so $\mathcal{F}_g$
and $\mathcal{F}_0$ are kept notationally separate precisely because ``ground'' and ``the
direction gravity points'' are not always the same thing, even though they agree at the launch
point.

$\mathcal{F}_r$ is right-handed, with $\hat{x}_r$ along the rocket's nose axis (positive forward),
$\hat{z}_r$ completing the pitch plane with it, and $\hat{y}_r$ out of the pitch plane. The pitch
angle $\theta_r$ is measured from $\hat{x}_0$ (the ground-frame horizontal), positive nose-up; the
pitch rate is $\omega_r=\mathrm{d}\theta_r/\mathrm{d}t$, the rotation rate about $\hat{y}_r$, and
the angular acceleration is $\varepsilon_r=\mathrm{d}\omega_r/\mathrm{d}t$. In ground-frame
components,
\begin{equation}\label{eq:body_axes_in_ground}
    \hat{x}_r = \cos\theta_r\,\hat{x}_0 - \sin\theta_r\,\hat{z}_0, \qquad
    \hat{z}_r = \sin\theta_r\,\hat{x}_0 + \cos\theta_r\,\hat{z}_0.
\end{equation}
($\alpha$ is reserved for the angle of attack, \autoref{sec:relwind}.)

$\mathcal{F}_t$ carries the actual line of thrust action, $\hat{x}_t$, which need not coincide
with the nose axis: a real nozzle can be canted, or otherwise misaligned, by a small angle
$\theta_t$ from the body axis --- a fixed property of how the motor is mounted, not of the flight
state. Giving thrust its own frame, rather than folding $\theta_t$ into every thrust term by hand,
is what lets that distinction be written down once:
\begin{equation}\label{eq:thrust_frame}
    \hat{x}_t = \cos\theta_t\,\hat{x}_r - \sin\theta_t\,\hat{z}_r, \qquad
    \hat{z}_t = \sin\theta_t\,\hat{x}_r + \cos\theta_t\,\hat{z}_r, \qquad \hat{y}_t\equiv\hat{y}_r.
\end{equation}
$\mathcal{F}_t$ is obtained from $\mathcal{F}_r$ exactly as $\mathcal{F}_r$ is from
$\mathcal{F}_0$ in \autoref{eq:body_axes_in_ground}, so a positive $\theta_t$ tilts the thrust line
in the same sense as a positive $\theta_r$ tilts the nose. The older cant angle $\delta$ had the
opposite sign, $\theta_t=-\delta$ (\autoref{sec:notation}). The thrust direction can in principle
be controlled; for now $\theta_t$ is carried as a constant input, ideally zero --- a fixed
manufacturing-type misalignment, swept by hand to see how sensitive the airframe is to it. A
commanded, time-varying gimbal angle driven by a control law is future work
(\autoref{sec:backlog}).

\subsection{Forces and moments}\label{sec:forces}
Weight is the main force the rocket has to overcome to increase its altitude:
\begin{equation}
    \vec{W} = mg\hat{z}_g,
\end{equation}
where $m$ is the rocket's mass and $g$ the gravitational acceleration along the local vertical
$\hat{z}_g$. Both factors are models rather than constants, and are chosen independently: the
magnitude $g$ and the direction $\hat{z}_g$ answer different questions --- how strong, and which
way --- and are supplied separately by \autoref{sec:gravity}.

Thrust is generated by ejecting propellant at high velocity, plus a contribution from the pressure
difference across the nozzle exit. With the thrust frame of \autoref{eq:thrust_frame} it is simply
\begin{equation}
    \vec{T} = T\hat{x}_t = T\left(\cos\theta_t\,\hat{x}_r - \sin\theta_t\,\hat{z}_r\right),
\end{equation}
purely axial in its own frame, with $T$ the thrust magnitude --- the measured thrust curve plus a
nozzle-exit pressure correction (\autoref{sec:thrust}); the angle $\theta_t$ inside $\hat{x}_t$
carries the whole of the misalignment, so the first form never needs to be rewritten as $\theta_t$
changes.

Motion through the surrounding air produces an aerodynamic force, and, about a reference point $p$
on the rocket (\autoref{sec:geometry}), an aerodynamic pitching moment:
\begin{equation}
    \vec{A}_f = -qS\left(C_x\hat{x}_r + C_z\hat{z}_r\right), \qquad
    \vec{A}_m = -qSc\,C_m\hat{y}_r, \qquad
    q = \tfrac12\rho v_{rel}^2,
\end{equation}
with $q$ the dynamic pressure, $\rho$ the air density, $v_{rel}$ the airspeed (the magnitude of
the rocket's velocity relative to the air, \autoref{sec:relwind}), $S$ the aerodynamic reference
area and $c$ the reference length. $C_x,C_z,C_m$ are the axial-force, normal-force and
pitching-moment coefficients, defined in the body frame (\autoref{sec:aero}). The aerodynamic
force and moment are \emph{always} these dynamic-pressure-scaled coefficient products, whichever
of \autoref{sec:aero}'s coefficient models is selected; no model in this document supplies an
aerodynamic force directly as a force.

\subsection{Geometry}\label{sec:geometry}
Each of the three forces acts at its own point on the rocket:
\begin{itemize}
    \item $A$ --- nose tip, the origin of every length below;
    \item $g$ --- centre of gravity, where weight acts;
    \item $t$ --- thrust application point, where the (possibly canted) thrust line of
    $\mathcal{F}_t$ crosses the body;
    \item $p$ --- aerodynamic reference point (centre of pressure), where $\vec{A}_f$ acts and
    about which $\vec{A}_m$ is taken.
\end{itemize}
Every dimension is defined in the rocket frame $\mathcal{F}_r$: $l$ is the axial distance of a
point from $A$, measured aft (positive towards the tail), and $r$ its radial distance from the
axis, along $-\hat{z}_r$. Writing $(a_1,a_2,a_3)_r\equiv a_1\hat{x}_r+a_2\hat{y}_r+a_3\hat{z}_r$
for components in $\mathcal{F}_r$,
\begin{equation}
    \vec{Ag} = -l_g\hat{x}_r - r_g\hat{z}_r, \qquad
    \vec{At} = -l_t\hat{x}_r - r_t\hat{z}_r, \qquad
    \vec{Ap} = -l_p\hat{x}_r.
\end{equation}
The moments below are taken about $g$, so what they need is the position of the other two points
relative to it:
\begin{equation}\label{eq:arms}
    \vec{gt} = -\vec{Ag} + \vec{At} = (l_g-l_t,\;0,\;r_g-r_t)_r, \qquad
    \vec{gp} = -\vec{Ag} + \vec{Ap} = (l_g-l_p,\;0,\;r_g)_r.
\end{equation}

The radial distances $r$ describe the axis misalignment: in a real rocket the centre of gravity
and the nozzle cannot be strictly on the axis, whether from manufacturing tolerance or from an
as-built mass asymmetry. $r$ can be zero, and is by default. As $p$ is a reference point rather
than a measured one, no misalignment is defined for it ($r_p\equiv0$); any asymmetry in the
external surfaces instead shows up in the aerodynamic coefficients themselves, e.g.\ in polar
plots, not in a radial offset of $p$. $r_g$ and $r_t$ are carried as ordinary constants alongside
$l_g,l_t$ throughout what follows, on the same footing as the thrust angle $\theta_t$ --- setting
them to zero recovers the on-axis geometry exactly, and there is no separate ``model'' for them
beyond that: like $\theta_t$, they are inputs representing an as-built offset, in principle
obtained from manufacturing metrology of the actual airframe rather than derived from anything
else in this document. $l_t$, by contrast, only matters when the thrust is canted: it enters the
moment through $\sin\theta_t$ alone (\autoref{eq:eom_rot_expanded}), so for $\theta_t=0$ its value
is immaterial and it can be set to zero.

\subsection{Equations of motion}\label{sec:eom}
\begin{gather}
    m\vec{a}_r = \sum \vec{F} = \vec{W} + \vec{T} + \vec{A}_f \label{eq:eom_trans}\\
    I\varepsilon_r = \sum \vec{M}(g) \cdot \hat{y}_r
    = \left[\vec{gt}\times\vec{T} + \vec{gp}\times\vec{A}_f + \vec{A}_m\right]\cdot \hat{y}_r
    \label{eq:eom_rot}
\end{gather}
Substituting the forces of \autoref{sec:forces},
\begin{equation}\label{eq:eom_trans_expanded}
    m\vec{a}_r = mg\hat{z}_g + T\left(\cos\theta_t\,\hat{x}_r - \sin\theta_t\,\hat{z}_r\right)
    - qS\left(C_x\hat{x}_r + C_z\hat{z}_r\right).
\end{equation}
For the rotation, substitute \autoref{eq:arms} and use the $\hat{y}_r$ component of the cross
product of two pitch-plane vectors,
$\left[(a_1,0,a_3)_r\times(b_1,0,b_3)_r\right]\cdot\hat{y}_r = a_3b_1-a_1b_3$:
\begin{align}
    I\varepsilon_r
    &= \left[T\left(-\vec{Ag}+\vec{At}\right)\times\hat{x}_t
       + \left(-\vec{Ag}+\vec{Ap}\right)\times\vec{A}_f + \vec{A}_m\right]\cdot\hat{y}_r
       \nonumber\\
    &= \Bigl[T\,(l_g-l_t,\,0,\,r_g-r_t)_r\times(\cos\theta_t,\,0,\,-\sin\theta_t)_r
       \nonumber\\
    &\qquad + qS\,(l_g-l_p,\,0,\,r_g)_r\times(-C_x,\,0,\,-C_z)_r
       - qSc\,C_m\hat{y}_r\Bigr]\cdot\hat{y}_r
       \nonumber\\
    &= T\left[(r_g-r_t)\cos\theta_t + (l_g-l_t)\sin\theta_t\right]
       + qS\left[-r_gC_x + (l_g-l_p)C_z - c\,C_m\right]. \label{eq:eom_rot_expanded}
\end{align}
Each term reads off the geometry directly. The radial offset $r_g-r_t$ of the thrust point gives
an arm to the axial thrust component $T\cos\theta_t$, and its axial offset $l_g-l_t$ an arm to the
normal component $T\sin\theta_t$; in the same way $r_g$ is the arm of the axial aerodynamic force
$qSC_x$ and $l_g-l_p$ that of the normal one $qSC_z$. For the aligned rocket,
$r_g=r_t=\theta_t=0$, thrust drops out of the moment altogether and
$I\varepsilon_r=qS\left[(l_g-l_p)C_z-c\,C_m\right]$.

Moments are taken about the centre of gravity $g$, so weight contributes none:
$\vec{gg}\times\vec{W}=0$, for whatever direction $\vec{W}$ points. $m$ and $I$ are not constants
once the propellant model of \autoref{sec:mass} is active --- they become $m(t), I(t)$, and the
centre of gravity $g$ itself travels as $l_g(t)$. Reinterpreting $\mathcal{F}_0$ as Earth-fixed
rather than inertial adds one further translational term to \autoref{eq:eom_trans} and nothing at
all to \autoref{eq:eom_rot} (\autoref{sec:noninertial}). \autoref{sec:analysis} expands these two
equations by hand.

These two equations describe \emph{free} flight, which is all of a flight except its two
mechanically constrained phases. While the rocket is still on the launch rail
(\autoref{sec:rail}) and while it hangs under a deployed recovery canopy (\autoref{sec:recovery}),
the attitude state is frozen, $\omega_r=\varepsilon_r=0$, and \autoref{eq:eom_trans} is modified
--- projected onto the rail direction in the first case, given one additional inertial-frame force
in the second. Both are written up with the models rather than here, because each needs its own
constraint or force model to be stated first; the equations above are what they modify, and what
they revert to.

\subsubsection{Model inventory}\label{sec:eom_inventory}
Every symbol on the right-hand sides of \autoref{eq:eom_trans_expanded} and
\autoref{eq:eom_rot_expanded} is either a state or is supplied by a model:
\begin{itemize}
    \item $m$, $I$ --- mass and moment of inertia, from the mass model (\autoref{sec:mass});
    \item $g\hat{z}_g$ --- gravity, magnitude and direction (\autoref{sec:gravity});
    \item $T$ --- thrust curve plus pressure correction (\autoref{sec:thrust});
    \item $\theta_t$ --- thrust direction: can be controlled, for now a constant, ideally zero
    (\autoref{sec:thrust});
    \item $q$ --- from the atmosphere model, the state and the wind model, through $\rho$ and
    $v_{rel}$ (\autoref{sec:atmosphere}, \autoref{sec:wind}, \autoref{sec:relwind});
    \item $S$, $c$ --- reference area and reference length (\autoref{sec:aero});
    \item $C_x$, $C_z$, $C_m$ --- aerodynamic coefficients (\autoref{sec:aero});
    \item $l_t$, $r_t$, $l_p$, $l_g$, $r_g$ --- lengths (\autoref{sec:geometry}); $l_g(t)$ moves
    with the mass model (\autoref{sec:mass}), and the position of $p$ is supplied together with the
    coefficients (\autoref{sec:aero});
    \item states --- $\theta_r$, $\omega_r$, $\varepsilon_r$ and $\vec{x}_r$, $\vec{v}_r$,
    $\vec{a}_r$ (\autoref{sec:frames}); the equations deliver $\varepsilon_r$ and $\vec{a}_r$, the
    rest is integrated;
    \item and, outside free flight, the parachute (\autoref{sec:recovery}) and the launch rail
    (\autoref{sec:rail}).
\end{itemize}

\subsection{Notation}\label{sec:notation}
The table maps each symbol of this report to the identifier the simulator uses for it
(\texttt{sim/} keeps its own names) and to the symbol used in earlier versions of this report.
``Not stored'' means the code keeps only an arm relative to the centre of gravity, not the length
from the nose tip; ``not implemented'' means the quantity appears in the report but not in the
simulator.
{\small
\begin{longtable}{>{\raggedright\arraybackslash}p{2.6cm}
                >{\raggedright\arraybackslash}p{2.0cm}
                >{\raggedright\arraybackslash}p{4.1cm}
                >{\raggedright\arraybackslash}p{4.4cm}
                >{\raggedright\arraybackslash}p{1.5cm}}
\hline
\textbf{Report} & \textbf{Earlier doc} & \textbf{Code} & \textbf{Meaning} & \textbf{Units} \\
\hline
\endfirsthead
\hline
\textbf{Report} & \textbf{Earlier doc} & \textbf{Code} & \textbf{Meaning} & \textbf{Units} \\
\hline
\endhead
\hline
\endfoot
\multicolumn{5}{l}{\textbf{Geometry and frames}} \\
\hline
$A$ & $A$ & --- & nose tip, origin of all lengths & --- \\
$g$ & $B$ & --- & centre of gravity & --- \\
$t$ & $C$ & --- & thrust application point & --- \\
$p$ & $D$ & --- & aerodynamic reference point (centre of pressure) & --- \\
$l_g$ & $l_B$ & \texttt{mass.cog\_position(m\_p, p)}; local \texttt{l\_B} in \texttt{sim/dynamics.py} (mass model only) & axial distance $A\to g$ & m \\
$r_g$ & $u_B$ & \texttt{p.lengths.u\_B} & radial offset of $g$ & m \\
$l_t$ & $l_C$ & not stored; $l_t-l_g$ = \texttt{p.lengths.l1 + p.lengths.l2} & axial distance $A\to t$ & m \\
$r_t$ & $u_C$ & \texttt{p.lengths.u\_C} & radial offset of $t$ & m \\
$l_p$ & $l_D$ & not stored; $l_p-l_g$ = \texttt{l1} (row below) & axial distance $A\to p$ & m \\
$r_p\equiv0$ & $u_D\equiv0$ & --- & no radial offset of $p$, by definition & m \\
$l_p-l_g$ & $l_1=l_D-l_B$ & \texttt{p.lengths.l1}; \texttt{aero.coefficients(\dots).l1} (varies under the table model) & arm $g\to p$, positive with $p$ aft & m \\
$l_t-l_p$ & $l_2=l_C-l_D$ & \texttt{p.lengths.l2} & arm $p\to t$ & m \\
$\vec{gt}$, $\vec{gp}$ & $\vec{BC}$, $\vec{BD}$ & \texttt{aB}, \texttt{aA} in \texttt{sim/dynamics.py} & position of $t$, $p$ relative to $g$ & m \\
$\theta_t$ & $-\delta$ & $-$\texttt{np.radians(}\allowbreak\texttt{p.lengths.misalign\_deg)} & thrust-line angle from $\hat{x}_r$ & rad (code: deg) \\
\hline
\multicolumn{5}{l}{\textbf{Mass}} \\
\hline
$m$, $I$ & $m$, $I$ & \texttt{p.weight.m}, \texttt{p.weight.I}, or \texttt{mass.moment\_of\_inertia} & mass; pitch moment of inertia about $g$ & kg, $\mathrm{kg\,m^2}$ \\
$m_r$ & $m_d$ less the casing & \texttt{airframe.m} (\texttt{configs/schema.py}); \texttt{p.weight.m\_d} also includes the casing & dry-rocket mass (constant) & kg \\
$m_m(t)$ & $m_p$ (propellant only) & \texttt{mass.propellant\_mass(m, p)} (\texttt{y[M]} $-$ \texttt{m\_d}) is the propellant only; $m_m$ adds the casing $m_f$ & motor mass (casing and propellant), burned down to $m_f$ & kg \\
$m_0$, $m_f$ & --- & \texttt{motor.m\_total}; $m_f$ = \texttt{m\_total - m\_prop} & motor mass at ignition, at burnout & kg \\
$\Delta t$ & --- & not implemented (no linear mass model) & burn duration (linear mass model) & s \\
$\Delta m$ & --- & \texttt{motor.m\_prop} & burned mass $m_0-m_f$ & kg \\
$\mathrm{ISP}$ & $I_{sp}$ & \texttt{p.weight.I\_sp}; \texttt{motor.I\_sp} & specific impulse & s \\
$g_0$ & $g_0$ & \texttt{p.g0} & standard gravity & $\mathrm{m/s^2}$ \\
$I_r$, $I_{r,g}$ & $I_d$ & \texttt{p.weight.I\_d} (dry assembly about its own centroid) & dry-rocket inertia about $g$; about its own centre of gravity & $\mathrm{kg\,m^2}$ \\
$I_m(t)$, $I_{m,g}(t)$ & $I_p(t)$ & \texttt{p.weight.I\_p\_spec}$\cdot m_p$ (propellant only) & motor inertia about $g$; about its own centre of gravity & $\mathrm{kg\,m^2}$ \\
$I_0$, $I_f$ & --- & \texttt{motor.I\_0}, \texttt{motor.I\_f} & $I_{m,g}$ at ignition, at burnout & $\mathrm{kg\,m^2}$ \\
$l_{g,r}$, $r_{g,r}$ & $l_d$ & \texttt{p.weight.l\_d}; radial: not stored (constant \texttt{u\_B}) & dry-rocket centre of gravity (axial, radial) & m \\
$l_{g,m}(t)$, $r_{g,m}(t)$ & $l_p$ (propellant) & \texttt{p.weight.l\_p} (fixed); radial: not stored & motor centre of gravity (axial, radial) & m \\
$l_0$, $l_f$ & --- & \texttt{motor.l\_cog\_0}, \texttt{motor.l\_cog\_f} & $l_{g,m}$ at ignition, at burnout & m \\
\hline
\multicolumn{5}{l}{\textbf{Forces and aerodynamics}} \\
\hline
$S$ & $S$ & \texttt{p.aero.S} & aerodynamic reference area (same symbol in both) & $\mathrm{m^2}$ \\
$c$ & $d_{ref}$ & \texttt{p.aero.d\_ref} & reference length of the pitching moment & m \\
$C_x$, $C_z$, $C_m$ & $C_x$, $C_z$, $C_m$ ($C_A$, $C_N$) & \texttt{p.aero.C\_x}, \texttt{C\_z}, \texttt{C\_m} (constant model); fields of \texttt{aero.AeroCoeffs} & axial-force, normal-force, pitching-moment coefficients & --- \\
$qSC_x$, $qSC_z$ & $A_x$, $A_z$ & \texttt{forces.aero\_force\_body} returns $(-A_x,-A_z)$ & aerodynamic force components & N \\
$qSc\,C_m$ & $A_m$ & \texttt{forces.aero\_moment\_y} returns $-A_m$ & aerodynamic pitching moment & N\,m \\
$S_k$, $c_k$, $C_{x,k}$, $C_{z,k}$, $C_{m,k}$, $pp_k$ & --- & not implemented & aero component $k$: reference area and length, coefficients, reference point & $\mathrm{m^2}$, m, ---, m \\
$C_{x,0}$, $k$ & $c_{x,0}$, $1/(e\pi\lambda)$ & not implemented & drag polar $C_x=C_{x,0}+kC_z^2$ & --- \\
$\partial C/\partial\alpha$, $\partial C/\partial M$, $\partial C/\partial Re$ & --- & only the $\alpha$ slope is in the code & linearized-model slopes & 1/rad, ---, --- \\
$C_{z,\alpha}$, $C_{m,\alpha}$, $C_{m,q}$, $C_{z,0}$, $C_{m,0}$ & $c_{z,\alpha}$, $c_{m,\alpha}$, $c_{m,q}$, \dots & aero coefficient inputs (\texttt{c\_mq} column of the aero table, \dots) & coefficient slopes and offsets & 1/rad or --- \\
\hline
\multicolumn{5}{l}{\textbf{Thrust}} \\
\hline
$T$, $\vec{T}$ & $T$, $\vec{T}$ & \texttt{thrust\_curve.thrust}; \texttt{forces.thrust\_body} & thrust magnitude; vector in $\mathcal{F}_r$ & N \\
$T_0(t)$ & $T_{ref}(t)$ & \texttt{thrust\_curve} (\texttt{T\_ref}) & reference (test-stand) thrust & N \\
$p_0$ & $p_{ref}$ & \texttt{p.thrust.p\_ref} & ambient pressure during the static test & Pa \\
$A_e$ & $A_e$ & \texttt{p.thrust.A\_e} & nozzle exit area & $\mathrm{m^2}$ \\
\hline
\multicolumn{5}{l}{\textbf{Atmosphere and wind}} \\
\hline
$p_{atm}$ & $p_{amb}$ & \texttt{air\_state(\dots).p} & ambient pressure & Pa \\
$q$ & $q$ & \texttt{flow.q} & dynamic pressure & Pa \\
$v_{rel}$, $\vec{v}_{rel}$ & $V$ & \texttt{flow.V} & airspeed; relative-wind vector & m/s \\
$v_{rel,x}$, $v_{rel,z}$ & $u$, $w$ & --- & relative-velocity components in $\mathcal{F}_r$ & m/s \\
$\alpha$ & $\alpha$ & \texttt{flow.alpha} & angle of attack & rad \\
$w_{up}$ & --- & not implemented & updraft speed ($w_z=-w_{up}$) & m/s \\
$v_{ref}$ & $V_{ref}$ & turbulence setting & fixed reference airspeed of the Dryden spectrum & m/s \\
\hline
\multicolumn{5}{l}{\textbf{States and integration}} \\
\hline
$\vec{y}=(x,z,v_x,$ $v_z,\theta_r,\omega_r,m)$ & $\vec{y}$ & \texttt{y[X,}\allowbreak\texttt{Z,}\allowbreak\texttt{VX,}\allowbreak\texttt{VZ,}\allowbreak\texttt{THETA,}\allowbreak\texttt{OMEGA,}\allowbreak\texttt{M]} & integrated state vector & m, m/s, rad, rad/s, kg \\
$\theta_r$ & $\theta$ & \texttt{y[THETA]} & pitch angle & rad \\
$\omega_r$ & $\omega$ & \texttt{y[OMEGA]} & pitch rate & rad/s \\
$\varepsilon_r$ & $\mathrm{d}\omega/\mathrm{d}t$ & \texttt{M\_y / I} ($\mathrm{d}\omega/\mathrm{d}t$ in \texttt{sim/dynamics.py}) & angular acceleration & $\mathrm{rad/s^2}$ \\
$\vec{x}_r$ & $\vec{x}=(x,z)$ & \texttt{y[X]}, \texttt{y[Z]} & position in $\mathcal{F}_0$ & m \\
$\vec{v}_r$ & $\mathrm{d}\vec{x}/\mathrm{d}t$ & \texttt{y[VX]}, \texttt{y[VZ]} & velocity in $\mathcal{F}_0$ & m/s \\
$\vec{v}_g$ & $\vec{v}_B$ & \texttt{y[VX]}, \texttt{y[VZ]} & CoG velocity (the same vector as $\vec{v}_r$) & m/s \\
$v_x$, $v_z$ & $\mathrm{d}x/\mathrm{d}t$, $\mathrm{d}z/\mathrm{d}t$ & \texttt{y[VX]}, \texttt{y[VZ]} & velocity components in $\mathcal{F}_0$ & m/s \\
$\vec{a}_r$ & $\mathrm{d}\vec{v}/\mathrm{d}t$ & \texttt{a\_i} in \texttt{sim/dynamics.py} & acceleration in $\mathcal{F}_0$ & $\mathrm{m/s^2}$ \\
$m_{launch}$ & --- & \texttt{s0.m}; \texttt{p.weight.m} & liftoff mass $m(0)$ & kg \\
$h$ & $dt$ & \texttt{p.solver.dt} & fixed RK4 step & s \\
$\theta_{rail}$, $L_{rail}$, $s_{rail}$, $\hat{e}_{rail}$ & same & \texttt{p.site.rail\_angle\_deg}, \texttt{p.site.L\_rail}, \texttt{rail\_distance}, \texttt{\_e\_rail} & rail elevation, length, distance travelled, direction & rad, m, m, --- \\
\hline
\multicolumn{5}{l}{\textbf{Recovery and air brake}} \\
\hline
$S_{ch}$, $C_{D,ch}$, $\vec{F}_{ch}$, $v_{term}$ & $S_c$, $C_{D,c}$, $\vec{F}_c$, $V_{term}$ & \texttt{p.recovery.S}, \texttt{p.recovery.Cd}, \texttt{forces.}\allowbreak\texttt{parachute\_}\allowbreak\texttt{drag\_}\allowbreak\texttt{inertial} & canopy area, drag coefficient, drag force, terminal speed & $\mathrm{m^2}$, ---, N, m/s \\
$F_b$ & $A_{brake}$ & \texttt{control.airbrake\_*} & air-brake axial force & N \\
\hline
\multicolumn{5}{l}{\textbf{Analysis}} \\
\hline
$\hat{\omega}=\omega_r c/(2v_{rel})$ & $\hat{q}$ & --- & non-dimensional pitch rate & --- \\
$M_{\theta_t}=(l_g-l_t)\,T\sin\theta_t$ & $M_\delta$ & --- & misalignment disturbance torque & N\,m \\
$s_i$ & $l_i$ & --- & position of a body point relative to the CoG ($C_{m,q}$ derivation) & m \\
$M_\alpha$, $M_q$, $\omega_n$, $\zeta$ & same & --- & short-period parameters & $\mathrm{1/s^2}$, 1/s, rad/s, --- \\
\end{longtable}
}
$A$ is only the nose tip, the reference area being $S$ ($A_e$, $A_f$, $A_m$ carry their own
subscripts). Some letters carry two meanings and are told apart by context only: $g$ is both the
centre of gravity and the gravitational acceleration, and as a subscript both the point ($l_g$,
$r_g$) and the gravity frame ($\mathcal{F}_g$, $\hat{z}_g$); $t$ as a point is likewise distinct
from time $t$, and $h$ is the RK4 step in \autoref{sec:states} but the altitude elsewhere. The
mass model of \autoref{sec:mass} uses the subscript $m$ for the motor and $r$ for the dry rocket,
so $l_p$ is only the aerodynamic reference point; in the code \texttt{p.weight.l\_p} is the
propellant centroid (the motor's $l_{g,m}$).

\end{aitext}

\newpage
\begin{aitext}
\section{Analytical analysis}\label{sec:analysis}
This section works through the pitching moment --- thrust on or off, thrust axis aligned or
misaligned --- and the static and dynamic stability results that follow. An acceleration component
such as $\vec{a}_r\cdot\hat{x}_r$ is the acceleration resolved along the body axes at the current
instant, not a derivative taken inside the rotating frame $\mathcal{F}_r$.

Two inputs are zero by default and reduce everything below to its simplest form: the thrust angle
$\theta_t$ and the gravity-direction rotation $\varphi=x/R_e$, which stays zero if the local
vertical is frozen at the launch point.

\subsection{Forces}\label{sec:analysis_forces}
Aerodynamic force and thrust are already body-axis quantities, so only weight has to be resolved.
Weight acts along the local vertical $\hat{z}_g(\varphi)=-\sin\varphi\,\hat{x}_0+\cos\varphi\,
\hat{z}_0$, so with \autoref{eq:body_axes_in_ground}
\begin{equation}
    \vec{W}\cdot\hat{x}_r = -mg\sin(\theta_r+\varphi), \qquad
    \vec{W}\cdot\hat{z}_r = +mg\cos(\theta_r+\varphi).
\end{equation}
The downrange rotation of gravity thus enters only as a shift of the pitch angle,
$\theta_r\to\theta_r+\varphi$: the pitch above the \emph{local} horizon. Adding thrust and the
aerodynamic force,
\begin{align}
    m\,\vec{a}_r\cdot\hat{x}_r &= -mg\sin(\theta_r+\varphi) + T\cos\theta_t - qSC_x \label{eq:analytic_force_x}\\
    m\,\vec{a}_r\cdot\hat{z}_r &= mg\cos(\theta_r+\varphi) - T\sin\theta_t - qSC_z \label{eq:analytic_force_z}
\end{align}
Sign checks, with $\varphi=\theta_t=0$. Nose straight up ($\theta_r=90^\circ$): weight is purely
axial, $m\,\vec{a}_r\cdot\hat{x}_r=-mg+T-qSC_x$. Nose horizontal ($\theta_r=0$): weight is purely
normal, $m\,\vec{a}_r\cdot\hat{z}_r=mg-qSC_z$. For $\varphi$, flying level at downrange $x>0$ gives
$m\,\vec{a}_r\cdot\hat{x}_r=-mg\sin\varphi<0$: the Earth's centre lies below and behind the
rocket, so the pull is tilted against the flight direction. At $x=10\,\mathrm{km}$ this is
$g\varphi\approx1.5\times10^{-2}\,\mathrm{m/s^2}$.

\subsection{Moment}\label{sec:analysis_moment}
Substituting the arms of \autoref{eq:arms} into \autoref{eq:eom_rot}, with $\hat{x}_r\times
\hat{z}_r=-\hat{y}_r$ and the thrust vector of \autoref{eq:thrust_frame}:
\begin{equation} \label{eq:moment_full}
    I\varepsilon_r = (r_g - r_t)T\cos\theta_t + (l_g - l_t)T\sin\theta_t
                     - r_g\,qSC_x + (l_g - l_p)\,qSC_z - qSc\,C_m
\end{equation}
Thrust acts at $t$, displaced from the CoG by $\vec{gt}=(l_g-l_t)\hat{x}_r+(r_g-r_t)\hat{z}_r$:
the radial offset gives an arm to the axial thrust component $T\cos\theta_t$, the axial offset an
arm to the normal component $T\sin\theta_t$. The aerodynamic force acts at $p$: the axial arm
$l_g-l_p$ carries $qSC_z$ and the radial arm $r_g$ carries $qSC_x$. Weight does not appear, since
$\vec{gg}\times\vec{W}=0$ for every orientation of $\vec{W}$. Rotating the local vertical is
therefore a purely translational correction, like the Coriolis term.

For on-axis thrust, $\vec{T}=T\hat{x}_r$, sliding the application point along the axis changes
nothing: the force is parallel to the displacement. Only $r_t$ survives. Once $\theta_t\neq0$ the
normal component does care how far along the axis it acts, hence the second term.

\autoref{eq:moment_full} specialises four ways; \emph{aligned} means $r_g=r_t=\theta_t=0$.
\begin{center}
\begin{tabular}{l l l}
 & \textbf{aligned} ($\theta_t=0$) & \textbf{misaligned} ($\theta_t\neq0$) \\
\textbf{no thrust} ($T=0$) & \autoref{eq:moment_nothrust} & identical --- $T=0$ removes $\theta_t$ \\
\textbf{with thrust} ($T>0$) & \autoref{eq:moment_aligned} & \autoref{eq:moment_misaligned} \\
\end{tabular}
\end{center}

\subsubsection{No thrust: coast and descent}\label{sec:analysis_nothrust}
Setting $T=0$ leaves the airframe's purely aerodynamic behaviour, independent of the motor:
\begin{equation} \label{eq:moment_nothrust}
    I\varepsilon_r = -r_g\,qSC_x + (l_g - l_p)\,qSC_z - qSc\,C_m.
\end{equation}
Since $T=0$ removes both thrust terms, misalignment is irrelevant here, and the static-stability
result below, built on this case, needs no revision when it is switched on.

\subsubsection{With thrust, aligned}\label{sec:analysis_aligned}
With $r_g=r_t=0$ and $\theta_t=0$,
\begin{equation} \label{eq:moment_aligned}
    I\varepsilon_r = (l_g - l_p)\,qSC_z - qSc\,C_m.
\end{equation}
Thrust has dropped out exactly: an aligned rocket's attitude dynamics are the same during boost
as during coast, however large $T$ is.

\subsubsection{With thrust, misaligned}\label{sec:analysis_misaligned}
With the nozzle canted by a small fixed $\theta_t$ and the geometry otherwise on-axis,
\begin{equation} \label{eq:moment_misaligned}
    I\varepsilon_r = \underbrace{(l_g - l_t)T\sin\theta_t}_{M_{\theta_t}}
                     + (l_g - l_p)\,qSC_z - qSc\,C_m.
\end{equation}
Unlike every aerodynamic term, $M_{\theta_t}$ survives $q=0$ and $\alpha=0$: while the motor burns,
the airframe carries a disturbance torque proportional to the thrust. Its sign: $l_t>l_g$ for any
real rocket, so $l_g-l_t<0$. A positive $\theta_t$ pushes the tail towards $-\hat{z}_r$ and pitches
the nose \emph{down}, $M_{\theta_t}<0$; a negative $\theta_t$ pitches it up.

Since $\partial M_{\theta_t}/\partial\alpha=0$, the stability derivative of
\autoref{eq:dtheta_dalpha} is untouched: misalignment shifts the equilibrium without changing
whether it is stable. Setting $\varepsilon_r=0$ with the linearised coefficients of
\autoref{eq:aero_linearised} and $C_{z,0}=C_{m,0}=0$ gives the trim angle of attack,
\begin{equation} \label{eq:alpha_trim}
    \alpha_{trim} = -\frac{M_{\theta_t}}{I\,\partial\varepsilon_r/\partial\alpha}
        = \frac{-(l_g - l_t)T\sin\theta_t}
               {qS\bigl[(l_g - l_p)C_{z,\alpha} - c\,C_{m,\alpha}\bigr]}.
\end{equation}
A stable airframe has a negative denominator, so $\alpha_{trim}$ carries the sign of
$M_{\theta_t}$: the rocket flies at a small constant angle off the relative wind, whose normal
force curves the trajectory, rather than tumbling. $\alpha_{trim}\propto T/q$ grows without bound
just after ignition and in any high-altitude burn, and as the static margin tends to zero --- the
physical reason launch rails exist.

The radial offsets $r_g,r_t$ enter \autoref{eq:moment_full} on the same footing as $\theta_t$,
adding the constant offset $(r_g-r_t)T\cos\theta_t-r_g\,qSC_x$; the worked cases set them to zero
to isolate $\theta_t$. A commanded, time-varying $\theta_t(t)$ remains future work.

\subsection{Static stability}\label{sec:analysis_stability}
Static stability asks whether a disturbance in the angle of attack $\alpha$ produces a restoring
moment. Take the no-thrust aligned case, \autoref{eq:moment_nothrust} with $r_g=0$, and linearise
the coefficients,
\begin{equation} \label{eq:aero_linearised}
    C_z = C_{z,0} + \alpha\,C_{z,\alpha}, \qquad
    C_m = C_{m,0} + \alpha\,C_{m,\alpha},
\end{equation}
with slopes $C_{z,\alpha}=\partial C_z/\partial\alpha$, $C_{m,\alpha}=\partial C_m/\partial\alpha$ and
offsets $C_{z,0}$, $C_{m,0}$. This is an analysis device applied to whichever coefficients are in
force; the Linearized model (\autoref{sec:aero_coeffs}) is its form in $\alpha$, $M$ and $Re$.
Stability requires $\partial\varepsilon_r/\partial\alpha<0$: a disturbance must produce an angular
acceleration that reduces it. The derivative is
\begin{equation} \label{eq:dtheta_dalpha}
    \frac{\partial \varepsilon_r}{\partial \alpha} = \frac{qS}{I}\biggl[(l_g - l_p)C_{z,\alpha} - c\,C_{m,\alpha}\biggr],
\end{equation}
negative for $C_{z,\alpha}>0$ (the normal case) when
\begin{equation} \label{eq:static_margin}
    l_g - l_p < c\,\frac{C_{m,\alpha}}{C_{z,\alpha}}.
\end{equation}
Both lengths are measured from the nose, so $l_g-l_p<0$ exactly when $p$ sits aft of $g$ --- the
familiar ``centre of pressure behind centre of gravity'' rule, which is what
\autoref{eq:static_margin} reduces to when $p$ is the centre of pressure itself, $C_{m,\alpha}=0$.
In this report's sign of $\vec{A}_m$, a positive $C_{m,\alpha}$ is stabilising, the opposite of
the textbook $C_{m,\alpha}<0$.

\paragraph{Static margin, in calibers.} $p$ is a moment reference station, not necessarily the centre
of pressure (\autoref{sec:geometry}): the CP is the derived station at which the moment slope about it
vanishes,
\begin{equation}
    l_{cp} = l_p - c\,\frac{C_{m,\alpha}}{C_{z,\alpha}},
\end{equation}
and the rocketry convention divides the CP-to-CoG arm by the reference length, flipping the sign so a
stable rocket gets a positive number:
\begin{equation}
    \mathrm{SM} = \frac{l_{cp}-l_g}{c} = \frac{l_p-l_g}{c} - \frac{C_{m,\alpha}}{C_{z,\alpha}}.
\end{equation}
$\mathrm{SM}>0$ is then equivalent to \autoref{eq:static_margin}; the two coincide, $l_{cp}=l_p$, exactly
when $C_{m,\alpha}=0$ about $p$ (the degenerate case where $p$ already is the CP).

\subsection{Flight-path angle and the point-mass gravity turn}\label{sec:analysis_flightpath}
Everything above tracks the \emph{nose} direction $\theta_r$. The trajectory itself is set by the
\emph{velocity} direction, and the two differ by the angle of attack. Define the flight-path angle
$\theta_v$, the angle of $\vec{v}_r$ above the horizon (conventionally written $\gamma$; here
$\theta_v$, since $\gamma_r$ is the roll angle of \autoref{sec:roll}), by the same construction as
\autoref{eq:body_axes_in_ground}:
\begin{equation}\label{eq:vel_axes_in_ground}
    \hat{x}_v = \cos\theta_v\,\hat{x}_0 - \sin\theta_v\,\hat{z}_0, \qquad
    \hat{z}_v = \sin\theta_v\,\hat{x}_0 + \cos\theta_v\,\hat{z}_0, \qquad
    \vec{v}_r = V\hat{x}_v,
\end{equation}
so that $v_x=V\cos\theta_v$, $v_z=-V\sin\theta_v$ and $\theta_v=\operatorname{atan2}(-v_z,v_x)$:
positive climbing, since $\hat{z}_0$ points down. In the pitch plane and in still air the three
angles close exactly,
\begin{equation}\label{eq:alpha_closure}
    \alpha = \theta_r - \theta_v,
\end{equation}
which is the simulator's $\alpha=\operatorname{atan2}(v_{rel,z},v_{rel,x})$ of
\autoref{sec:relwind} written in ground-frame terms. With wind, $\theta_v$ in
\autoref{eq:alpha_closure} is the air-relative flight-path angle.

\paragraph{The point-mass limit.} Strip the airframe of its rotational degree of freedom: no
aerodynamic normal force, thrust along the velocity vector ($\theta_t=0$ and $\alpha=0$, so
$\hat{x}_v=\hat{x}_r$), a flat non-rotating Earth ($\hat{z}_g=\hat{z}_0$, $\varphi=0$), still air
(so the path speed is the airspeed, $V=v_{rel}$) and drag $D$
antiparallel to $\vec{v}_r$. Then \autoref{eq:eom_trans} reduces to
$m\vec{a}_r = mg\hat{z}_0 + (T-D)\hat{x}_v$, while differentiating
\autoref{eq:vel_axes_in_ground} gives
\begin{equation}\label{eq:accel_in_vel_axes}
    \vec{a}_r = \dot{V}\hat{x}_v - V\dot\theta_v\,\hat{z}_v .
\end{equation}
Projecting on $\hat{z}_v$, where $\hat{z}_0\cdot\hat{z}_v=\cos\theta_v$ and weight is the only
contributor, and on $\hat{x}_v$:
\begin{gather}
    \dot\theta_v = -\frac{g}{V}\cos\theta_v, \label{eq:flightpath_rate}\\
    m\dot{V} = T - D - mg\sin\theta_v. \label{eq:speed_rate}
\end{gather}
\autoref{eq:flightpath_rate} is the \emph{gravity turn}: gravity alone bends the velocity vector, at
a rate $g/V$ modulated by $\cos\theta_v$. It vanishes when the flight is vertical
($\theta_v=\pm90^\circ$), where gravity is purely decelerating and turns nothing, and is largest
when level, $\dot\theta_v=-g/V$. A faster rocket turns less per unit time --- the reason a high
thrust-to-weight launch holds its initial attitude. Holding $V$ fixed, \autoref{eq:flightpath_rate}
integrates to
\begin{equation}\label{eq:gravity_turn_closed}
    \tan\!\left(\frac{\pi}{4}+\frac{\theta_v}{2}\right)
        = \tan\!\left(\frac{\pi}{4}+\frac{\theta_{v,0}}{2}\right) e^{-gt/V},
\end{equation}
so $\theta_v\to-90^\circ$ with time constant $V/g$.

This is a \emph{cruder} model than \autoref{eq:eom_trans}--\autoref{eq:eom_rot}, not a refinement of
them: it carries no moment, no inertia and no angle of attack, and the simulator never integrates it.
Its uses are the two below --- a closed-form estimate to check the integrator against, and the
justification of the short-period approximation that follows.

\subsubsection{Dynamic stability}\label{sec:analysis_dynamic}
The static criterion does not say what the resulting oscillation does. That needs the
\textbf{pitch-damping derivative} $C_{m,q}$, the sensitivity of $C_m$ to the non-dimensional pitch
rate $\hat{\omega}=\omega_r c/(2v_{rel})$:
\begin{equation}
    C_m = C_{m,0} + \alpha\,C_{m,\alpha} + \hat{\omega}\,C_{m,q}.
\end{equation}
It comes from the body's own rotation. A point a distance $s_i$ aft of the CoG moves relative to
the CoG with velocity $\omega_r\hat{y}_r\times\vec{s}_i$ and sees a local angle of attack shifted
by $\Delta\alpha_i\approx+\omega_r s_i/v_{rel}$. Fins well aft see the largest shift, and the
resulting normal-force change against its arm gives a moment proportional to $\omega_r$, present
even at $\alpha=0$.

With the sign of $\vec{A}_m$ used here, a positive $C_m$ pitches the nose \emph{down}. The simulator's
code flags this non-standard sign by renaming the internal quantity to \texttt{C\_m\_code}, keeping the
plain name \texttt{C\_m} reserved at the code's I/O boundary (config files, CSV aero-table headers) for
the standard convention when a user declares \texttt{cm\_convention: standard}. A statically stable,
aft-finned airframe therefore has $C_{m,q}>0$: a positive (nose-up) pitch rate
produces an additional positive (nose-down) moment, opposing the rotation.

Linearising \autoref{eq:moment_aligned} with this term (and $C_{z,0}=C_{m,0}=0$),
\begin{equation}
    I\varepsilon_r = qS(l_g-l_p)C_{z,\alpha}\,\alpha - qSc\,C_{m,\alpha}\,\alpha
        - \frac{qSc^2}{2v_{rel}}C_{m,q}\,\omega_r,
\end{equation}
which separates into a stiffness term and a damping term,
\begin{equation}
    M_\alpha \equiv \frac{\partial\varepsilon_r}{\partial\alpha}
        = \frac{qS}{I}\bigl[(l_g-l_p)C_{z,\alpha}-c\,C_{m,\alpha}\bigr], \qquad
    M_q \equiv \frac{qSc^2}{2v_{rel}I}C_{m,q}.
\end{equation}
$M_\alpha$ is negative for a statically stable airframe and $M_q$ is positive for $C_{m,q}>0$.
Eliminating $\theta_r$ in favour of $\alpha$ through \autoref{eq:alpha_closure}
($\omega_r=\dot\alpha+\dot\theta_v$, $\varepsilon_r=\ddot\alpha+\ddot\theta_v$) turns
$\varepsilon_r=M_\alpha\alpha-M_q\omega_r$ into an exact statement,
\begin{equation}\label{eq:short_period_exact}
    \ddot\alpha + M_q\dot\alpha - M_\alpha\alpha
        = -\left(\ddot\theta_v + M_q\dot\theta_v\right),
\end{equation}
whose right-hand side is the entire content of the short-period approximation: dropping it gives
\begin{equation}
    \frac{\mathrm{d}^2\alpha}{\mathrm{d}t^2} + M_q\frac{\mathrm{d}\alpha}{\mathrm{d}t} - M_\alpha\alpha = 0,
\end{equation}
with natural frequency and damping ratio
\begin{equation}
    \omega_n = \sqrt{-M_\alpha}, \qquad \zeta = \frac{M_q}{2\omega_n} = \frac{M_q}{2\sqrt{-M_\alpha}},
\end{equation}
valid only for $M_\alpha<0$; otherwise the disturbance grows rather than oscillates.

\autoref{eq:flightpath_rate} makes the dropped term quantitative instead of merely ``slow''. With
$|\dot\theta_v|\le g/V$ and $\ddot\theta_v=O\!\left((g/V)^2+g\dot V/V^2\right)$, the dominant
neglected term is $M_q\dot\theta_v=2\zeta\omega_n\dot\theta_v$. Measuring it against the retained
stiffness $M_\alpha\alpha\sim\omega_n^2\hat\alpha$ for an oscillation of amplitude $\hat\alpha$, and
writing
\begin{equation}\label{eq:flightpath_ratio}
    \Lambda \equiv \frac{g/V}{\omega_n} = \frac{g}{V\sqrt{-M_\alpha}},
\end{equation}
the approximation holds when $2\zeta\Lambda\ll\hat\alpha$ and $\Lambda^2\ll\hat\alpha$ --- in short,
$\Lambda\ll1$: the flight-path angle must drift through much less than the oscillation amplitude in
one short period, $\Delta\theta_v\approx2\pi\Lambda\ll\hat\alpha$. Since $\Lambda\propto1/V$, the
approximation is at its weakest at low airspeed --- on and just off the rail, and near the apogee of
a lofted flight --- and strongest through the high-$q$ part of the boost. Near-vertical flight
relaxes it independently of speed, because $\cos\theta_v\to0$ suppresses $\dot\theta_v$ directly.

The distributed-force argument above gives an estimate. A component $i$ with normal-force slope
$C_{z,\alpha,i}$ (referred to $S$) at a distance $s_i$ from the CoG adds a moment proportional to
$\omega_r s_i^2/v_{rel}$, and equating the sum with the $C_{m,q}$ term gives
\begin{equation}
    C_{m,q} \approx \frac{2}{c^2}\sum_i C_{z,\alpha,i}\,s_i^2 ,
\end{equation}
positive in this report's sign convention. OpenRocket uses a different estimate, from the cross-flow
drag of the rotating body and fins \cite{ortech}.
\restodo{a value of $C_{m,q}$ for the project's airframe: the component slopes $C_{z,\alpha,i}$, or
a measured value (forced-oscillation tests or unsteady CFD)}

\end{aitext}

\newpage
\begin{aitext}
\section{Models}\label{sec:models}
Each subsection gives the interchangeable models behind one term of the equations of motion: the
three forces (gravity, thrust with its mass, aerodynamics), the environment they depend on
(atmosphere, wind, relative flow), the Earth-rotation correction, and the integration. The last
four subsections --- launch rail, recovery, air brake and sampled-data mode --- modify the
equations rather than supply a term, and are off by default.

\subsection{Gravity}\label{sec:gravity}
This subsection supplies the weight $\vec{W}=mg\hat{z}_g$: the magnitude $g$ and the direction
$\hat{z}_g$ (down, along the local vertical), chosen independently. Below, $h$ is the altitude above
the surface and $R_e$ the Earth's mean radius.

\subsubsection{Magnitude}
\begin{itemize}
    \item \textbf{Constant:} $g=g_0$.
    \item \textbf{Radial} (inverse-square, spherical non-rotating Earth).
    \item \textbf{Ellipsoidal} (WGS84, experimental): Somigliana surface gravity plus a
    second-order free-air correction, with $\phi$ the geodetic latitude of the launch site and
    constants from \autoref{sec:constants}.
\end{itemize}
\begin{gather}
    g(h)=g_0\left(\frac{R_e}{R_e+h}\right)^2 \qquad \text{(radial)} \\
    \gamma_0(\phi) = \gamma_e\frac{1+k\sin^2\phi}{\sqrt{1-e^2\sin^2\phi}} \\
    g(\phi, h) = \gamma_0\left[1 - \frac{2}{a_e}\left(1+f+m-2f\sin^2\phi\right)h + \frac{3}{a_e^2}h^2\right].
\end{gather}
The ellipsoidal magnitude already contains the centrifugal part of the Earth's rotation.

\subsubsection{Direction}
\begin{itemize}
    \item \textbf{Parallel} (default): $\hat{z}_g\equiv\hat{z}_0$, the local vertical at the launch
    point, for the whole flight.
    \item \textbf{Radial:} $\hat{z}_g$ points to the centre $O$ of a sphere of radius $R_e$.
    \item \textbf{Ellipsoidal:} $\hat{z}_g$ is the inward normal of the WGS84 ellipsoid.
\end{itemize}
With $\vec{Og}=\vec{OO^*}+\vec{O^*g}$, $\vec{OO^*}=-R_e\hat{z}_0$ and $\vec{O^*g}=x\hat{x}_0+z\hat{z}_0$
in the pitch plane,
\begin{align}
    \text{parallel:}\quad & \hat{z}_g = \hat{z}_0, \\
    \text{radial:}\quad & \hat{z}_g = -\frac{\vec{Og}}{|\vec{Og}|}
        = -\sin\varphi\,\hat{x}_0 + \cos\varphi\,\hat{z}_0, \qquad \tan\varphi=\frac{x}{R_e-z}.
        \label{eq:zhat_g_rotated}
\end{align}
The code uses the arc angle $\varphi=x/R_e$ (\autoref{eq:phi_downrange}), which agrees with the
exact angle at $z=0$ and to first order elsewhere.

For the ellipsoid the normal does not pass through $O$: it makes the angle $\phi-\psi$ with $\vec{Og}$,
where $\psi$ is the geocentric latitude, $\tan\psi=(1-e^2)\tan\phi$ (at most $0.19^\circ$). The
normal at the rocket is the normal at $O^*$ turned by the change of geodetic latitude along the
track, $\vartheta=\int_0^x \mathrm{d}x'/R_n(x')$, with $R_n$ the local normal-curvature radius, so
\begin{equation}\label{eq:zhat_g_ellipsoid}
    \hat{z}_g = -\sin\vartheta\,\hat{x}_0 + \cos\vartheta\,\hat{z}_0 .
\end{equation}
The code freezes $R_n$ at its launch value, the radius of the osculating sphere in the flight
azimuth $\chi$ (Euler's formula, with $M,N$ the meridional and prime-vertical radii of curvature),
\begin{equation}
    R_\chi = \left(\frac{\cos^2\chi}{M(\phi_0)} + \frac{\sin^2\chi}{N(\phi_0)}\right)^{-1},
    \quad M = \frac{a_e(1-e^2)}{(1-e^2\sin^2\phi)^{3/2}}, \quad
    N = \frac{a_e}{\sqrt{1-e^2\sin^2\phi}},
\end{equation}
so $\vartheta=x/R_\chi$, i.e.\ \autoref{eq:zhat_g_rotated} with $R_e\to R_\chi$.
\todo{this differs from the exact normal at third order in $\vartheta$, because $R_n$ keeps
changing with latitude along the track.}

\paragraph{Position relative to $O$, direction and true altitude.} The non-parallel models need
three things: where the rocket is relative to $O$, the resulting $\hat{z}_g$ (above), and the true
altitude $h_{true}$ over the curved ground that the altitude-dependent models use in place of
$-z$. The arc angle and the ground point under the rocket are
\begin{gather}
    \varphi = \frac{x}{R_e}, \label{eq:phi_downrange}\\
    \vec{O^*P}(\varphi) = R_e\sin\varphi\,\hat{x}_0 + R_e(1-\cos\varphi)\,\hat{z}_0,
        \label{eq:surface_point}
\end{gather}
and the altitude is the offset of the rocket from that ground point along $-\hat{z}_g$,
\begin{equation} \label{eq:h_true}
    h_{true}(x,z) = \left[x\hat{x}_0+z\hat{z}_0-\vec{O^*P}(\varphi)\right]\cdot\bigl(-\hat{z}_g\bigr)
        = x\sin\varphi - z\cos\varphi - R_e(1-\cos\varphi).
\end{equation}
For small $\varphi$, $h_{true}\approx x^2/2R_e-z$. Impact is $h_{true}=0$ instead of $z=0$. For the
ellipsoid, $R_e\to R_\chi$.

\subsubsection{Vector field (magnitude + direction)}
A single model may supply $g$ and $\hat{z}_g$ together, as a vector field
$\vec{g}(\vec{x},t)$ from a gridded dataset, so that $\vec{W}=m\vec{g}$.
\todo{no vector-field gravity model exists in sim/; only the separate magnitude and direction
models are implemented.}

\subsection{Thrust}\label{sec:thrust}
\subsubsection{Magnitude}
Magnitude of thrust is the result of conservation of momentum (ejection of fuel) and pressure
differences, which can be computed analytically for an internal-combustion nozzle:
\begin{equation}
    T = \left|\frac{\mathrm{d}m_m}{\mathrm{d}t}\right| v_e + A_e(p_e - p_{atm}),
\end{equation}
with the motor mass flow, the exhaust speed $v_e$, the nozzle-exit pressure
$p_e$, the exit area $A_e$ and the ambient pressure $p_{atm}$.
In real life, however, the data from a static test is used instead. The \emph{reference} thrust
curve $T_0(t)$ is given as scatter (data) points $(t_i,T_i)$, $i=1,\dots,N$, read from a motor file,
and evaluated by linear interpolation between neighbouring points, held at $T_1$ for $t<t_1$ and
zero for $t>t_N$; burnout is the falling zero-crossing of this curve:
\begin{equation}
    T_0(t) = \mathrm{lerp}\!\left(t; \{t_i\}, \{T_i\}\right).
\end{equation}
As the test is static, the ambient pressure $p_0$ it was measured at is fixed and does not follow
the rocket's altitude. To account for the change in ambient pressure $p_{atm}$ during flight, the
net thrust once the motor is flowing is
\begin{equation}\label{eq:thrust_pressure}
    T(t, h) = T_0(t) + A_e\left(p_0 - p_{atm}(h)\right),
\end{equation}
valid only while $T_0(t)>0$ (i.e.\ while the motor is actually flowing); $A_e$ is the nozzle exit area and $p_0$ the ambient pressure the
curve was measured at (the test-stand pressure). Since $p_{atm}$ falls with altitude, the thrust rises as the rocket climbs, by at most
$A_ep_0$ in vacuum. $A_e=0$ recovers $T_0(t)$ exactly.
\todo{sim/ names $p_0$ \texttt{p.thrust.p\_ref} and $A_e$ \texttt{p.thrust.A\_e}; the
default $A_e=0$ switches the correction off.}

\subsubsection{Direction}
The line of thrust need not coincide with the nose axis: a real nozzle can be canted, or otherwise
misaligned, by a small angle $\theta_t$ from the body axis. The unit vector along the line of
thrust, $\hat{x}_t$, is obtained from $\hat{x}_r$ exactly as $\hat{x}_r$ is from $\hat{x}_0$:
\begin{equation}\label{eq:thrust_frame}
    \hat{x}_t = \cos\theta_t\,\hat{x}_r - \sin\theta_t\,\hat{z}_r, \qquad
    \hat{z}_t = \sin\theta_t\,\hat{x}_r + \cos\theta_t\,\hat{z}_r, \qquad \hat{y}_t\equiv\hat{y}_r,
\end{equation}
so a positive $\theta_t$ tilts the thrust line in the same sense as a positive $\theta_r$ tilts the
nose. Three options:
\begin{itemize}
    \item \textbf{Ideal} (the default): $\theta_t=0$, the thrust acts along $\hat{x}_r$.
    \item \textbf{Misalignment}: $\theta_t$ is a constant input --- a manufacturing-type error,
    chosen once per airframe, not a function of $t$ or $h$.
    \item \textbf{Controlled orientation} (gimbal): $\theta_t(t)$ commanded by a control law, the
    only pitch authority a rocket has once it is off the rail. Future work.
\end{itemize}
$\theta_t$ does not scale the magnitude, and the mass flow is driven by $T_0(t)$ regardless of
where the nozzle points. The axial component delivered is $T\cos\theta_t$ (a $4\times10^{-5}$
loss at $0.5^\circ$, negligible), while $T\sin\theta_t$ acts at the arm $l_t-l_g$ and is not
negligible as a moment.

\subsection{Mass, centre of gravity and moment of inertia}\label{sec:mass}
As propellant burns, mass leaves the rocket through the nozzle and both the centre of gravity and
the moment of inertia change over the burn. The rocket is therefore split into two bodies: the
\emph{dry rocket} (subscript $r$: everything that does not burn, with constant mass $m_r$, centre
of gravity $l_{g,r}$ from the nose tip $A$, and moment of inertia $I_{r,g}$ about its own centre of
gravity) and the \emph{motor} (subscript $m$: casing and remaining propellant, with time-dependent
mass $m_{m}(t)$, centre of gravity $l_{g,m}(t)$ and moment of inertia $I_{m,g}(t)$ about its own
centre of gravity). The subscript $g$ on $I_{r,g}$, $I_{m,g}$ says ``about that body's own centre
of gravity''; the total quantities $m$, $l_g$, $I$ carry no body subscript. $t$ is time since
ignition. (The subscript $p$ is not used for the propellant: it names only the aerodynamic
reference point.) $g_0$ is the standard gravity, a fixed constant of the definition of specific
impulse and unrelated to the local $g(h)$.

\subsubsection{Mass}\label{sec:mass_mass}
The total mass is the dry rocket plus the motor,
\begin{equation}\label{eq:mass_total}
    m(t) = m_r + m_m(t).
\end{equation}
Only $m_m(t)$ changes; its rate is fixed by one of two motor mass models.

\paragraph{Linear model.} Inputs: the motor mass at ignition $m_0$, at burnout $m_f$ (the casing and
whatever else stays; $m_f=0$ if $m_m$ counts propellant only), and the burn duration $\Delta t$.
The mass falls at a constant rate,
\begin{equation}\label{eq:mass_linear}
    \frac{\mathrm{d}m_m}{\mathrm{d}t} = \frac{m_f-m_0}{\Delta t}, \qquad
    m_m(t)=m_0+\frac{m_f-m_0}{\Delta t}\,t \quad(0\le t\le\Delta t),
\end{equation}
and $m_m=m_f$ afterwards. The rate is negative, since $m_f<m_0$. It ignores the shape of the
thrust curve: mass leaves at the same rate whether the motor is delivering its peak or its tail-off.

\paragraph{Specific-impulse model.} Inputs: the thrust curve $T_0(t)$ and the specific impulse
$\mathrm{ISP}$ (seconds). The propellant leaves at the rate that delivers exactly that thrust,
\begin{equation}\label{eq:mass_isp}
    \frac{\mathrm{d}m_m}{\mathrm{d}t} = -\frac{T_0(t)}{\mathrm{ISP}\;g_0},
\end{equation}
integrated alongside the equations of motion until $m_m$ reaches $m_f$. The minus sign is
essential: the motor mass decreases. $T_0(t)$ is the
\emph{reference} thrust curve (the test-stand curve), not the
pressure-corrected $T=T_0(t)+A_e(p_0-p_{atm})$. The reason is that $\mathrm{ISP}$ is
defined from the stand measurement: it is the ratio of the stand thrust to the propellant weight
flow, and the nozzle-exit pressure term is a change in the force delivered to the airframe, not in
the rate at which the propellant is consumed. When the mass model is built from the thrust file
alone, $\mathrm{ISP}$ follows from the curve's total impulse $J=\int T_0\,\mathrm{d}t$ and the
burned mass $\Delta m=m_0-m_f$ as $\mathrm{ISP}=J/(\Delta m\,g_0)$, so that the integral of
\autoref{eq:mass_isp} burns exactly $\Delta m$.
\todo{sim/dynamics.py computes the mass flow from \texttt{thrust\_curve.thrust}, i.e.\ the
pressure-corrected $T$, not from the bare $T_0(t)$; the two differ only when $A_e\neq0$ (default
$A_e=0$, so no current result changes). Decide whether the code should be switched to $T_0$, as
derived here.}

\paragraph{In the code.} Only the specific-impulse model is implemented, selected by
\texttt{motor.mass\_model} with the values \texttt{"none"}, \texttt{"isp"} and \texttt{"endpoints"};
the last obtains $\mathrm{ISP}$ from $J/(\Delta m\,g_0)$ as above
(\texttt{mass.isp\_from\_impulse}). There is no separate linear-in-time model of
\autoref{eq:mass_linear}; $m(t)$ is the integrated state \texttt{y[M]}.
\todo{linear mass model (\autoref{eq:mass_linear}) is not implemented in sim/.}

\subsubsection{Centre of gravity}\label{sec:mass_cog}
The first moment of mass about the nose tip $A$ is additive:
\begin{equation}\label{eq:mass_cog}
    l_g(t)\,m(t) = l_{g,r}\,m_r + l_{g,m}(t)\,m_m(t)
    \quad\Longrightarrow\quad
    l_g(t) = \frac{l_{g,r}\,m_r + l_{g,m}(t)\,m_m(t)}{m(t)},
\end{equation}
the CoG distance from the nose. The same
holds for the radial offset when it is not zero,
\begin{equation}\label{eq:mass_cog_r}
    r_g(t)\,m(t) = r_{g,r}\,m_r + r_{g,m}(t)\,m_m(t).
\end{equation}
Part of the motor's centre of gravity travel is the motor itself shifting aft as the forward
grain burns away. The simplest model of it is linear in the \emph{burned} mass: with $l_0$ and $l_f$
the motor's centre of gravity at ignition and burnout,
\begin{equation}\label{eq:mass_l_linear}
    \frac{\partial l_{g,m}}{\partial m_b}=\frac{l_f-l_0}{\Delta m},\qquad
    l_{g,m}(m_m)=l_0+\frac{l_f-l_0}{\Delta m}\,(m_0-m_m),
\end{equation}
where $m_b=m_0-m_m$ is the burned mass and $\Delta m=m_0-m_f$ its total. (The derivative is taken with respect to the burned mass $m_b$: with respect to $m_m$ itself the
sign would flip.) At $m_m=m_0$ this gives $l_0$, at $m_m=m_f$ it gives $l_f$.

\subsubsection{Moment of inertia}\label{sec:mass_inertia}
The moments of inertia add once each body's own moment is transferred to the total centre of
gravity $g$ by the parallel-axis theorem,
\begin{equation}\label{eq:mass_inertia}
    I(t) = I_r + I_m(t),
\end{equation}
\begin{equation}\label{eq:mass_inertia_parts}
    I_r = I_{r,g} + m_r\,d_r^2, \qquad I_m(t) = I_{m,g}(t) + m_m(t)\,d_m(t)^2,
\end{equation}
where $I_{r,g}$, $I_{m,g}(t)$ are taken about each body's own centre of gravity. In the
two-dimensional pitch plane the transfer distances are
\begin{equation}\label{eq:mass_dist}
    d_r^2=\bigl(l_g(t)-l_{g,r}\bigr)^2+\bigl(r_g(t)-r_{g,r}\bigr)^2, \qquad
    d_m^2=\bigl(l_g(t)-l_{g,m}(t)\bigr)^2+\bigl(r_g(t)-r_{g,m}(t)\bigr)^2.
\end{equation}
With all bodies on the axis ($r\equiv0$) the radial terms vanish, and \autoref{eq:mass_dist}
reduces to the axial form $d_r=l_g(t)-l_{g,r}$, $d_m=l_g(t)-l_{g,m}(t)$. Note that
$I_r$ depends on time through $l_g(t)$ although $I_{r,g}$ and $m_r$ do not.

The motor's own moment of inertia is modelled, like its centre of gravity, as linear in burned
mass: with $I_0$, $I_f$ its values at ignition and burnout,
\begin{equation}\label{eq:mass_I_linear}
    \frac{\partial I_{m,g}}{\partial m_b}=\frac{I_f-I_0}{\Delta m},\qquad
    I_{m,g}(m_m)=I_0+\frac{I_f-I_0}{\Delta m}\,(m_0-m_m),
\end{equation}
with $I_{m,g}=I_0$ at $m_m=m_0$ and $I_f$ at $m_m=m_f$.

\paragraph{Radial offsets.}
\autoref{eq:mass_cog_r} and \autoref{eq:mass_dist} are the general forms. The rest of this document
treats $r_g$ as a constant as-built input (zero by default), and the code does the same: the mass model is purely axial (\texttt{mass.cog\_position},
\texttt{mass.moment\_of\_inertia} take no radial argument) and \texttt{p.lengths.u\_B} is a constant.
For $r_{g,r}=r_{g,m}=0$ the two agree; for a rocket with an off-axis dry-rocket CoG, $r_g(t)$ would
drift as the motor burns, and the transfer distances would gain the radial terms, neither of which
is implemented.
\todo{radial offsets are inconsistent between the document and the code: \autoref{eq:mass_cog_r} and
\autoref{eq:mass_dist} make $r_g(t)$ time-dependent and add a radial term to the parallel-axis
distance, whereas sim/ keeps $r_g$ (\texttt{u\_B}) constant and omits the radial term. Exact for
$r\equiv0$, which is the default.}

\paragraph{Variable-mass terms.} Differentiating the angular momentum of a body whose mass and inertia
change gives, besides $I\varepsilon_r$, the terms $\omega_r\,\mathrm{d}I/\mathrm{d}t$ and the
jet-damping moment of the exhaust leaving at distance $l_t-l_g$ from the centre of gravity (of the
order of $|\mathrm{d}m_m/\mathrm{d}t|\,(l_t-l_g)^2\,\omega_r$). The equation of motion used here is
simply $I(t)\,\varepsilon_r=\sum M$; both terms are dropped.
\todo{variable-mass terms $\omega_r\,\mathrm{d}I/\mathrm{d}t$ and jet damping are neglected (sim/
integrates $I(t)\,\varepsilon_r=\sum M$ with no such terms); estimate their size against the
pitch damping for a short, high-thrust motor.}

\paragraph{In the code.} The motor is represented not by the linear models above but by the
two-lump composition of \texttt{configs/schema.py}: a casing lump and a propellant lump of mass
$m_{prop}$ with a \emph{fixed} propellant centroid and $I_{prop}\propto$ remaining propellant mass.
Equivalent to \autoref{eq:mass_cog}--\autoref{eq:mass_inertia_parts} with the dry body
the dry body of mass $m_r+m_f$ (airframe plus casing; \texttt{Weight.m\_d}, \texttt{l\_d},
\texttt{I\_d}) and the motor reduced to the propellant alone (\texttt{l\_p}, \texttt{I\_p\_spec}). The endpoint inputs
$l_0,l_f,I_0,I_f$ (\texttt{l\_cog\_0}, \texttt{l\_cog\_f}, \texttt{I\_0}, \texttt{I\_f}) are not
used to interpolate linearly: they are inverted at assembly time into the fixed-centroid
propellant lump, which reproduces $l_0,l_f,I_0,I_f$ at both ends but is not linear in between.
\todo{linear centre-of-gravity and moment-of-inertia models (\autoref{eq:mass_l_linear},
\autoref{eq:mass_I_linear}) differ from the two-lump model of sim/ in between the endpoints; they
agree only at ignition and burnout.}

\subsubsection{Propellant grain geometry}
The propellant part of $m_m(t)$, $l_{g,m}(t)$ and $I_{m,g}(t)$ depends on the grain geometry and on
how it recedes as it burns, so it deserves elaborating rather than just flagging. Solid-motor grains commonly come in a few
standard families: \textbf{cylindrical} grains (a solid or centrally bored cylinder, burning on
the outer surface, or --- the common BATES configuration --- on the inner bore and, if
uninhibited, the end faces too); \textbf{tubular/annular} grains (a thin-walled, case-bonded tube
burning on the inner bore only); \textbf{end-burning} grains (burning only on one flat face, which
gives a nearly constant burn area and hence a nearly constant thrust, but a centroid that moves
along the grain quite differently from a radially-burning one); and \textbf{star or finocyl}
grains (a star- or wagon-wheel-shaped central perforation, chosen specifically to tailor the
burn-area-versus-time curve, and so the thrust curve $T_0(t)$ itself). Each family gives a
different relationship between burn area and web thickness $A_b(w)$; combined with a regression
rate $\mathrm{d}w/\mathrm{d}t=r_b(p)$ (typically Saint-Robert's law, $r_b=ap^n$, with $p$ the
chamber pressure), the instantaneous propellant volume, its centroid, and its own moment of
inertia about that centroid all follow by direct integration over the receding grain shape; added
to the casing's contribution they give $l_{g,m}(t)$ and $I_{m,g}(t)$ of
\autoref{eq:mass_cog}--\autoref{eq:mass_inertia_parts}.

That is exactly the information the linear models of \autoref{eq:mass_l_linear} and
\autoref{eq:mass_I_linear} (and, in the code, the fixed-centroid propellant lump with
$I_{prop}\propto$ propellant mass) are standing in for: a centroid fixed or moving linearly and an
inertia proportional to the remaining mass assume a uniform density and a shape that merely shrinks in
overall size, self-similarly, as it burns --- which no real grain does. The true centroid and the
true inertia ratio both drift over the burn, and by how much depends on which family above is
used and on the specific grain's own dimensions: bore and outer diameter, grain length, and, for
star grains, the number and shape of the points. None of that is fixed anywhere in this document
yet, so the placeholder stays in place until an actual grain geometry is chosen; this section can
be revisited once one is.
\todo{change of mass or moment of inertia due to mechanisms (e.g.\ deployable surfaces) is not
covered above}
\todo{estimate the dry-rocket inputs $m_r$, $l_{g,r}$, $I_{r,g}$ from components (OpenRocket-style)
instead of entering them by hand.}

\subsection{Aerodynamics}\label{sec:aero}
\subsubsection{Force and moment}
The aerodynamic force and the pitching moment about the reference point $p$ are
\begin{equation} \label{eq:aero_magnitudes}
    \vec{A}_f = -qS\left(C_x\hat{x}_r + C_z\hat{z}_r\right), \qquad
    \vec{A}_m = -qSc\,C_m\hat{y}_r,
\end{equation}
with $q$ the dynamic pressure, $S$ the reference area, $c$ the reference length, and
$C_x,C_z,C_m$ the axial-force, normal-force and pitching-moment coefficients. The component
magnitudes are $qSC_x$, $qSC_z$ and $qScC_m$.

\paragraph{Reference dimensions.} $S$ is the area of the pipe, i.e.\ the cross-section of the body
tube, and $c$ is the diameter of the pipe; $c$ also forms the Reynolds number.

These equations are \emph{definitions}, not a model: they hold for every coefficient model, and
the force is never supplied directly. One pair $(S,c)$ per component is shared by all of its
coefficient models, since a coefficient means nothing except relative to the $(S,c)$ it was
normalised by.

\paragraph{Multi-component aerodynamics.} When the airframe is built from several aerodynamic
components $k$ (body, fins, nose cone, \dots), each with its own reference point $pp_k$, its own
force $\vec{A}_{f,k}$ and moment $\vec{A}_{m,k}$ about $pp_k$, the total force is the sum and the
total moment about the common reference point $p$ adds the moment of each force about $p$:
\begin{equation} \label{eq:aero_multi}
    \vec{A}_f = \sum_k \vec{A}_{f,k}, \qquad
    \vec{A}_m(p) = \sum_k \vec{A}_{m,k} + \sum_k \vec{p\,pp_k}\times\vec{A}_{f,k},
\end{equation}
where $\vec{p\,pp_k}$ is the position of $pp_k$ relative to $p$ (the lever arm of component $k$'s
force about $p$). Each component's coefficients are normalised by its \emph{own} reference area and length,
$\vec{A}_{f,k}=-qS_k(C_{x,k}\hat{x}_r+C_{z,k}\hat{z}_r)$ and $\vec{A}_{m,k}=-qS_kc_k\,C_{m,k}\hat{y}_r$,
so $S_k,c_k$ generally differ between components and from the airframe's $S,c$.
\todo{sim/ carries a single aerodynamic component (one $S,c$ and one coefficient set per airframe);
\autoref{eq:aero_multi} is not implemented.}

$l_p-l_g$ is the axial distance from the centre of gravity $g$ to the aerodynamic reference point
$p$, positive with $p$ aft, so that $\vec{gp}=-(l_p-l_g)\hat{x}_r$ and
the caliber static margin is $\mathrm{SM}=(l_p-l_g)/c$.

\paragraph{Axial force and the flow direction.} The axial term $-qSC_x\hat{x}_r$ of
\autoref{eq:aero_magnitudes} assumes nose-first flight, with the relative flow along $-\hat{x}_r$
and $|\alpha|$ small. The coefficient $C_x$ is defined in body axes and is positive, so the term
always points aft; when the rocket flies tail-first (after apogee on a ballistic descent, or in a
tumble) the relative flow arrives from behind, and the same aft-pointing term would \emph{push}
the rocket along its motion --- ``drag'' that accelerates it. The physically correct axial force
opposes the axial component of the relative flow: its sign must follow the sign of
$\hat{x}_r\cdot\vec{v}_{rel}$ (equivalently of $\cos\alpha$), not stay fixed.
\todo{\autoref{eq:aero_magnitudes} assumes nose-first flight. The correction (a sign factor on
$C_x$ that follows $\cos\alpha$, continuous through $\alpha=90^\circ$ and equal to $1$ for
nose-first flight) is a separate code decision; once accepted, state the factor here.}

\subsubsection{Dynamic pressure}
\begin{equation}
    q = \tfrac12\rho v_{rel}^2,
\end{equation}
with $\rho$ the local air density and $v_{rel}$ the magnitude of the rocket's velocity relative to
the air. Both $v_{rel}$ and the angle of attack $\alpha$ that feeds the coefficients below come
from the relative-flow model.

\subsubsection{Coefficients}\label{sec:aero_coeffs}
$C_x,C_z,C_m$ and the arm $l_p-l_g$ are not derived in this document from the airframe's shape ---
they are supplied. Three interchangeable models supply them, selected per run; all three feed the
definitions above. Writing $C$ for
any one of $C_x,C_z,C_m$, they are, in order of increasing fidelity: constant,
\begin{equation}
    C = C_0,
\end{equation}
linearized about a reference condition in the angle of attack, the Mach number and the Reynolds
number,
\begin{equation} \label{eq:aero_general_linear}
    C(\alpha,M,Re) = C_0 + \alpha\frac{\partial C}{\partial\alpha} + M\frac{\partial C}{\partial M}
        + Re\frac{\partial C}{\partial Re},
\end{equation}
and real data, a multivariable lookup table $C(\alpha,M,Re)$. The bullets below give the form each
takes in the code.
\todo{sim/ implements the $\alpha$ slope only in the linearized model (no $M$ or $Re$ slope), and
the table model is a two-variable $(M,\alpha)$ table; $Re$ enters neither.}
\begin{itemize}
    \item \textbf{Constant} (the cheapest model): $C_x,C_z,C_m$ are fixed inputs and
    $l_p-l_g$ is fixed, independent of both $\alpha$ and $M$. Useful for first-cut sizing, and for
    isolating one effect at a time in a comparison sweep --- freezing the aerodynamic coefficients
    is the cleanest way to attribute a difference between two runs to the model actually under
    study. Note what is \emph{not} frozen: the forces themselves still vary over the flight, since
    the constant coefficients are multiplied by $q=\tfrac12\rho v_{rel}^2$ like everybody else's.
    \item \textbf{Linearized}: a single linear slope in $\alpha$ about the zero-$\alpha$ condition,
    with the moment arm still fixed,
    \begin{equation} \label{eq:aero_linearized}
        C_x = C_{x,0}, \qquad C_z = C_{z,\alpha}\alpha, \qquad C_m = C_{m,\alpha}\alpha, \qquad l_p-l_g \text{ fixed}.
    \end{equation}
    This is \autoref{eq:aero_general_linear} with $C_{z,0}=C_{m,0}=0$ and no $M$ or $Re$ slope.
    $M$-independent, and valid only for small $\alpha$ --- which an on-trajectory flight of a
    statically stable airframe mostly is. Implementation note: $C_z,C_m$ are evaluated as
    $C_{z,\alpha}\sin\alpha$, $C_{m,\alpha}\sin\alpha$ in code (identical to the formula above for
    small $\alpha$, within $0.4\%$ at $8.6^\circ$) rather than the raw linear form. This is a
    numerical-robustness fix, not a different physical model:
    $\alpha=\operatorname{atan2}(w,u)$ is discontinuous at its own branch cut ($u<0$, $w$ crossing
    zero), which an unperturbed flight coasting tail-first through $v_{rel}=0$ at apogee lands on
    exactly, and the unbounded linear form turns that into a solver-breaking jump in the moment
    every evaluation. $\sin\alpha$ removes the discontinuity (and, at large $\alpha$, is arguably
    more correct besides: it is exactly zero at $\alpha=\pi$, matching the zero net side force a
    body of revolution has flying exactly tail-first by symmetry).
    \item \textbf{Table} (the highest-fidelity model available here): $C_x,C_z,C_m$ and the arm
    $l_p-l_g$ are read from a database, populated once per airframe from wind-tunnel testing, CFD
    sweeps, or a semi-empirical panel-method tool, and interpolated at run time. In practice this
    means a rectangular table over $(M,\alpha)$: bracket and clamp $M$ at the grid edges (no
    extrapolation beyond the table's Mach range), interpolate linearly in $\alpha$ within each
    bracketing row, then blend the two rows linearly across the $M$ bracket. This is the only model
    in which $l_p-l_g$, and hence the static margin, responds to the flight state at all.
\end{itemize}

\paragraph{Table file format.} The table is a CSV file \texttt{data/aero/<name>.csv}, read by
\texttt{sim/aero.py}. Lines starting with \texttt{\#} and blank lines are comments; the first other
line is the header, and every following line is one grid point $(M,\alpha)$:
\begin{center}
\small
\begin{tabular}{lll}
\hline
\textbf{Column} & \textbf{Meaning} & \textbf{Units} \\
\hline
\texttt{alpha\_deg} & angle of attack $\alpha$ (axis) & deg \\
\texttt{M} & Mach number (axis) & --- \\
\texttt{C\_A} & axial-force coefficient ($C_x$) & --- \\
\texttt{C\_N} & normal-force coefficient ($C_z$) & --- \\
\texttt{C\_m} & pitching-moment coefficient & --- \\
\texttt{l1\_m} & arm $l_p-l_g$ from the centre of gravity to $p$, positive aft & m \\
\texttt{c\_mq} & pitch-damping derivative $C_{m,q}$ (optional column) & --- \\
\hline
\end{tabular}
\end{center}
The grid must be rectangular: every $M$ value repeats exactly the same set of \texttt{alpha\_deg}
values, otherwise the loader raises an error. Interpolation is as described above (clamped in $M$,
linear in $\alpha$, then linear blend across the $M$ bracket); nothing is extrapolated.
\todo{the table has no Reynolds-number axis; a third axis $Re$ would extend the format by one
\texttt{Re} column and a trilinear interpolation. A long CSV (one row per grid point) is kept as the
exchange format since it is what CFD post-processing scripts export most easily.}

\paragraph{Pitch damping across the models.} The pitch-damping derivative $C_{m,q}$ ---
positive $C_{m,q}$ being damping in this document's sign of $\vec{A}_m$ --- is a property of the
airframe and applies under all three models. Constant and Linearized take it as a single constant
input; Table takes it from an optional column and falls back to that constant. It enters as an
additive contribution to $C_m$, the further term $\hat{\omega}\,C_{m,q}$ with
$\hat{\omega}=\omega_r c/(2v_{rel})$ the non-dimensional pitch rate. The moment contribution is
evaluated in the equivalent product form
\begin{equation} \label{eq:pitch_damping_moment}
    qSc\,\hat{\omega}\,C_{m,q} = \tfrac14\rho v_{rel}Sc^2\,\omega_r\,C_{m,q},
\end{equation}
which never forms $\hat{\omega}$ and so never divides by $v_{rel}$, which is zero at apogee.

The force is resolved along the body axes directly, since $C_x,C_z$ are body-axis coefficients.

The Table model's database has to cover the angle of attack $\alpha$, the Mach number $M$, the
Reynolds number $Re$ and the geometry (fin planform and cant, body-fin interference, surface
finish).
\todo{expand once the CFD/data-extraction pipeline that populates this database is defined; a
control surface (airbrake, aileron, etc.) has a defined point of entry --- a control law
perturbing $C_x,C_z,C_m$ or the arm $l_p-l_g$ --- and the implementation carries that
entry point as an identity-by-default seam, with no non-identity variant registered, but neither
the surface's geometry nor any deflection law is written down here; the air brake deliberately
does not use it, entering instead as a separate additive drag area; pitch rate is
covered above through $C_{m,q}$ as a constant or a table column, not as a function of $(\alpha,M)$
in its own right; and, for 3D, other stuff}

\subsubsection{Linearization of coefficients}\label{sec:aero_linearization}
Coefficients can be linearised in $\alpha$ for the static-stability analysis:
\begin{equation}
    C_z = C_{z,0} + \alpha\,C_{z,\alpha},
\end{equation}
\begin{equation}
    C_m = C_{m,0} + \alpha\,C_{m,\alpha},
\end{equation}
where $C_{z,\alpha}=\partial C_z/\partial\alpha$ and $C_{m,\alpha}=\partial C_m/\partial\alpha$ are the
slopes and $C_{z,0}$, $C_{m,0}$ the offsets. This is an analysis device applied to whichever
coefficients are in force; the \textbf{Linearized} model (\autoref{eq:aero_linearized}) is the same
two equations with $C_{z,0}=C_{m,0}=0$ and a constant $C_x=C_{x,0}$.

\paragraph{Optional drag polar.} Instead of supplying $C_x$ as an independent family, it can be
tied to the normal-force coefficient through an induced-drag term,
\begin{equation} \label{eq:drag_polar}
    C_x = C_{x,0} + k\,C_z^2,
\end{equation}
with $C_{x,0}$ the zero-lift axial coefficient and $k$ the induced-drag factor. For a lifting
surface (airfoil, wing) the closed form is $k=1/(e\pi\lambda)$, with $e$ the Oswald efficiency
factor and $\lambda$ the aspect ratio; for a body of revolution with fins no such closed form is
known here, so $k$ would be fitted to CFD or wind-tunnel data. The benefit is one coefficient
family less to tabulate, and an axial force that stays consistent with the normal force at an angle
of attack, instead of being tabulated separately and possibly inconsistently. The limits: it is an
approximation valid for attached flow (small to moderate $\alpha$); it is parabolic in $C_z$, so it
misses flow separation and the Mach-number dependence of $k$ unless $C_{x,0}$ and $k$ are
themselves tabulated in $M$; and the axial force on a slender body at $\alpha$ is not purely
induced drag.
\todo{decide whether to adopt the drag polar (\autoref{eq:drag_polar}). sim/ does not have it:
$C_x$ is constant, $C_{x,0}$ or a table column.}

\subsection{Atmosphere}\label{sec:atmosphere}
Goal: given altitude $h$ (and, for the interpolated model, horizontal position), provide the local
air density $\rho$, together with everything needed to form the Reynolds and Mach numbers of the
relative flow, and the local speed of sound $a$ and dynamic viscosity $\mu$.

\subsubsection{Model tiers}\label{sec:atmos_tiers}
Three tiers, increasing in fidelity, all exposing the same $(\rho, T, p, a, \mu)$ interface:
\begin{enumerate}
    \item \textbf{ISA} (International Standard Atmosphere): the standard piecewise-linear
    temperature-lapse model (troposphere + lower stratosphere), pressure from the hydrostatic
    equation, density from the ideal gas law, speed of sound $a=\sqrt{\gamma R_{specific}T}$,
    viscosity from Sutherland's law.
    \item \textbf{Real data} (one sounding): a profile ($T(h)$, $p(h)$) measured at one
    geographic location at one time, or an average profile for that location (e.g.\ a seasonal
    mean), interpolated in altitude only.
    \item \textbf{Real data across multiple stations}: the same kind of data from several
    geographic locations, interpolated in altitude, horizontal position, and optionally time, to
    estimate local conditions along the actual ground track. Implemented for the spatial
    interpolation (the \emph{Field} paragraph below); \todo{interpolation in time between
    soundings is not implemented -- each station contributes one profile.}
\end{enumerate}

\paragraph{Disturbances and gusts.} The three tiers above give the \emph{mean} state of the
air. Disturbances about that mean --- turbulence and gusts --- perturb the velocity of the air
and not its thermodynamic state, so they belong to the wind model. A
\emph{discrete gust} is a single deterministic event of prescribed shape (e.g.\ the ``1-minus-cosine''
ramp); \emph{continuous turbulence} is a stationary random process described by a spectrum. Only the
latter is modelled (Dryden). \todo{A
discrete-gust model is not in the code; add it as a deterministic wind increment $\Delta\vec{v}_w(t)$
if gust-response cases are wanted.}

\paragraph{ISA.} Geopotential altitude from geometric altitude above mean sea level $h_{msl}$:
\begin{equation}
    H = \frac{R_{isa}h_{msl}}{R_{isa}+h_{msl}}, \qquad R_{isa} = 6356766\,\mathrm{m}.
\end{equation}
Seven layers, base geopotential altitude $H_b$, lapse rate $L$, base temperature/pressure
$T_b, p_b$:
\begin{center}
\begin{tabular}{l c c c c}
Layer & $H_b$ [m] & $L$ [K/m] & $T_b$ [K] & $p_b$ [Pa] \\
troposphere  & 0     & $-0.0065$ & 288.15 & 101325   \\
tropopause   & 11000 & $0$       & 216.65 & 22632.06 \\
stratosphere & 20000 & $+0.0010$ & 216.65 & 5474.889 \\
stratosphere & 32000 & $+0.0028$ & 228.65 & 868.0187 \\
stratopause  & 47000 & $0$       & 270.65 & 110.9063 \\
mesosphere   & 51000 & $-0.0028$ & 270.65 & 66.93887 \\
mesosphere   & 71000 & $-0.0020$ & 214.65 & 3.956420 \\
\end{tabular}
\end{center}
Within the layer containing $H$, temperature is linear in geopotential altitude:
\begin{equation}
    T = T_b + L(H-H_b).
\end{equation}
For a layer with non-zero lapse rate the pressure follows the polytropic relation,
\begin{equation}
    p = p_b(T/T_b)^{-g_{isa}/(LR_d)},
\end{equation}
and for an isothermal layer ($L=0$) the exponential form instead,
\begin{equation}
    p = p_b\exp\!\left(\frac{-g_{isa}(H-H_b)}{R_dT_b}\right).
\end{equation}
Density, speed of sound and viscosity then follow from the ideal gas law and Sutherland's law:
\begin{equation}
    \rho = \frac{p}{R_dT}, \qquad a = \sqrt{\gamma R_dT}, \qquad
    \mu = \mu_0\left(\frac{T}{T_0}\right)^{3/2}\frac{T_0+S_\mu}{T+S_\mu},
\end{equation}
with $R_d = 287.0528\,\mathrm{J/(kg\,K)}$, $\gamma = 1.4$, $g_{isa} = 9.80665\,\mathrm{m/s^2}$ (the
ISA's own fixed value), $\mu_0 = 1.716\times10^{-5}\,\mathrm{Pa\,s}$,
$T_0 = 273.15\,\mathrm{K}$, $S_\mu = 110.4\,\mathrm{K}$.

\paragraph{Sounding.} One station's measured $(h_{msl}, p, T, u, v)$ profile, temperature and wind
interpolated linearly, pressure log-linearly:
\begin{equation}
    T(h) = \mathrm{lerp}(h;\{h_i\},\{T_i\}), \quad
    p(h) = \exp\mathrm{lerp}(h;\{h_i\},\{\ln p_i\}), \quad
    u(h),v(h) = \mathrm{lerp}(h;\{h_i\},\{u_i,v_i\}).
\end{equation}
$\rho, a, \mu$ are derived from the interpolated $T, p$ exactly as ISA above --- never interpolated
directly, to stay thermodynamically consistent. Below the lowest level the profile continues with
the sounding's own local lapse rate; above the highest, an ISA-shaped profile re-anchored at
$(h_{top}, T_{top}, p_{top})$ keeps $T, p$ continuous at the seam. Wind is held at the edge value
outside the data range.

\paragraph{Field.} $N$ stations projected onto a tangent plane centred at the launch site
$(\phi_0,\lambda_0)$, together with the rocket's own ground position $(x\sin\chi, x\cos\chi)$
($x$ = downrange component of $\vec{x}_r$, $\chi$ = launch heading, degrees clockwise from North, the direction
of $+\hat{x}_0$):
\begin{equation}
    E = R_e\cos\phi_0\,(\lambda-\lambda_0), \qquad N = R_e\,(\phi-\phi_0).
\end{equation}
The $N$ profiles are combined, weighted by inverse squared distance, into one profile:
\begin{equation}
    w_i = \frac{1}{d_i^2+\varepsilon^2}, \qquad \hat{w}_i = \frac{w_i}{\sum_j w_j}, \qquad
    (T, \ln p, u, v) = \sum_i \hat{w}_i\,(T_i, \ln p_i, u_i, v_i),
\end{equation}
then interpolated in altitude exactly as the sounding tier. A single station reduces every
$\hat{w}_i$ to 1 and reproduces the sounding tier exactly.

Inverse-distance weighting is the default horizontal combination, and is chosen for being defined
\emph{everywhere} --- including far outside the polygon the stations span, which a rocket's ground
track routinely leaves --- for reproducing a station's own profile at that station, for working
with one or two stations as happily as with $N$, and for never oscillating. $\varepsilon$ (a few
hundred metres) removes the singularity at zero distance. One alternative is available in its
place, a thin-plate-spline radial basis function over the same tangent-plane station coordinates,
smoother when there are many well-spread stations; it can overshoot between them, and it needs at
least three non-collinear stations, so it is offered as an alternative rather than a replacement.
Methods defined only inside the convex hull of the stations are deliberately not offered at all:
they return nothing at exactly the downrange distances this document cares about.

\subsubsection{Density}\label{sec:density}
Water vapour is lighter than dry air --- $M_v=18.02$ against $M_d=28.96\,\mathrm{g/mol}$ --- so at
a given pressure and temperature, replacing some of the air with vapour makes the mixture
\emph{less} dense, not more. Damp air is light air. The effect is at the percent level, and it is
one of the few corrections in this document that can be taken straight from measured data, because
radiosondes report dewpoint alongside temperature and pressure.

\paragraph{Partial pressures.} Treat the mixture as two ideal gases sharing a volume, dry air at
partial pressure $p-e$ and vapour at partial pressure $e$:
\begin{equation}
    \rho = \frac{p-e}{R_dT} + \frac{e}{R_vT}
         = \frac{p}{R_dT}\left[1 - \frac{e}{p}\left(1 - \epsilon\right)\right], \qquad
    \epsilon = \frac{R_d}{R_v} = \frac{287.0528}{461.5} = 0.6220.
\end{equation}
The bracket is the whole correction, and it is always less than one. Folding it into an effective
temperature instead of an effective density is the standard bookkeeping --- the \emph{virtual
temperature}, the temperature dry air would need in order to have the moist mixture's density at
the same pressure:
\begin{equation} \label{eq:virtual_temperature}
    T_v = \frac{T}{1 - \dfrac{e}{p}\left(1-\epsilon\right)}, \qquad
    \rho = \frac{p}{R_dT_v},
\end{equation}
so that every formula keeps the dry-air gas constant $R_d$ and only the temperature fed to it
changes. $T_v \geq T$ always, with equality for perfectly dry air.

\paragraph{Vapour pressure from dewpoint.} The dewpoint $T_d$ is \emph{defined} as the temperature
to which the parcel must be cooled, at constant pressure and mixing ratio, to reach saturation ---
so the actual vapour pressure of the parcel equals the \emph{saturation} vapour pressure evaluated
at the dewpoint, $e=e_s(T_d)$. That identity is what makes the dewpoint column directly usable.
For $e_s$ the Magnus form is used, with the Alduchov--Eskridge coefficients, $T_d$ in
$^\circ\mathrm{C}$ and $e_s$ in hPa:
\begin{equation} \label{eq:magnus}
    e_s(T_d) = 6.1094\,\exp\!\left(\frac{17.625\,T_d}{T_d + 243.04}\right).
\end{equation}
This form is accurate to better than $0.4\%$ over $-40$ to $+50\,^\circ\mathrm{C}$, which covers
every dewpoint a sounding reports. The factor of 100 converting hPa to Pa before forming $e/p$ is
the one easy mistake here.

\paragraph{Scope: measured tiers only.} The correction is an opt-in switch, off by default --- a
run that does not ask for it gets dry air, $\rho=p/(R_dT)$, exactly as before this subsection
existed. When it is on, it applies to the \textbf{sounding} and \textbf{field} tiers, and
\emph{not} to ISA. That is not an omission: the International Standard
Atmosphere is defined as dry air of fixed composition, and its pressure column is computed on that
basis --- injecting moisture into it would give a profile that is neither the standard nor a
measurement. A sounding with no dewpoint data falls back to
$\rho=p/(R_dT)$ level by level, so the correction is present exactly where there is data to
support it.

\paragraph{Magnitude, and what is deliberately left dry.} At a warm, humid surface
($T=30\,^\circ\mathrm{C}$, $T_d=26\,^\circ\mathrm{C}$, $p=101325\,\mathrm{Pa}$) the correction is
$-1.25\%$ in density; at a temperate surface ($T=14.2\,^\circ\mathrm{C}$, $T_d=10\,^\circ\mathrm{C}$)
$-0.48\%$; at a cold, dry one ($T_d=-33\,^\circ\mathrm{C}$) $-0.017\%$; near the tropopause it is
below $10^{-5}$ and vanishes on its own without needing a cutoff, because $e_s$ collapses with
temperature. So this is a boundary-layer effect in practice --- which is where dynamic pressure is
largest during the early ascent, so it is not negligible for that reason. The speed of sound and
the viscosity are deliberately \emph{not} corrected: a consistent moist $a$ needs the mixture's own
$\gamma$ as well as its $R$, the two effects partly cancel, and what remains is $\sim0.2\%$ in $a$,
far below the uncertainty of the Mach-dependent coefficients it would feed. The hydrostatic interpolation of $p$ between levels likewise stays on dry
$T$; strictly it should use $T_v$, but it is a within-layer interpolation of measured endpoints,
not an integration over a deep column.

\subsubsection{Reynolds and Mach numbers}\label{sec:reynolds_mach}
\begin{equation}
    Re = \frac{\rho v_{rel} c}{\mu}, \qquad M = \frac{v_{rel}}{a},
\end{equation}
with $v_{rel}$ the airspeed, $c$ the reference length (the body diameter), $\mu$ the dynamic
viscosity, and $a$ the local speed of sound, both from the tier above.

\subsection{Wind}\label{sec:wind}
Goal: given time $t$ and position (altitude $h$, and horizontal position for the interpolated
tier), provide the wind velocity $\vec{v}_w$ in $\mathcal{F}_0$, used in
$\vec{v}_{rel} = \vec{v}_r - \vec{v}_w$. Same tier structure as the atmosphere:
\begin{enumerate}
    \item \textbf{Simple analytic model}: e.g.\ constant wind plus an altitude-dependent shear
    profile, with an optional, separately selected turbulence term (below).
    \item \textbf{Real data} (one sounding): measured/forecast wind profile $\vec{v}_w(h)$
    at one location and time, or an average profile for that location.
    \item \textbf{Real data across multiple stations}: wind profiles from multiple locations,
    interpolated in altitude and horizontal position (inverse-distance weighting, as for the
    atmosphere); implemented in \texttt{sim/} \texttt{atmosphere.py}.
    \todo{Interpolation in time between soundings is not implemented.}
\end{enumerate}
The wind enters only through the relative velocity $\vec{v}_{rel}=\vec{v}_r-\vec{v}_w$; the
aerodynamic force and moment see the air through $\vec{v}_{rel}$ alone.

\paragraph{Flight-plane projection.} The launch heading $\chi$ fixes the 2D flight plane. A horizontal wind given as ENU components $(u,v)$, or as meteorological
speed/direction $(W,\psi)$ with $u=-W\sin\psi,\ v=-W\cos\psi$ ($\psi$ = the direction the wind
blows \emph{from}), projects onto it as
\begin{equation}
    w_x = u\sin\chi + v\cos\chi = -W\cos(\psi-\chi), \qquad
    w_y = u\cos\chi - v\sin\chi = -W\sin(\psi-\chi),
\end{equation}
with $w_x$ the downrange (along-plane) component feeding $\vec{v}_w$ and $w_z=0$ (no vertical wind
modelled). $w_y$, the cross-plane component, is computed and discarded: crosswind weathercocking
needs the out-of-plane freedom only a full 6DOF model has.

\paragraph{Analytic.} Power-law shear and constant direction:
\begin{equation} \label{eq:wind_shear}
    W(h) = W_{ref}\left(\max(h,h_{ref})/h_{ref}\right)^{n}, \qquad \psi = \psi_0.
\end{equation}
$h$ is altitude above the local ground, $n\approx1/7$ the shear exponent (the classic open-terrain
value; smoother terrain is lower, rougher/urban/forested terrain higher, up to $\sim0.4$), valid
within the atmospheric surface/boundary layer only --- roughly the lowest few hundred metres to
1--2\,km depending on stability, above which the sounding tiers below should be used instead. This
is the \emph{mean} profile only; the fluctuation about it is a separate, independently selected
model, below.

\paragraph{Sounding / field.} The same measured $(u,v)$ profiles and interpolation as
the atmosphere, projected through the formula above; never interpolated as $(W,\psi)$,
since direction wraps at $360^\circ$.

\paragraph{Turbulence.} A fluctuation added to the analytic tier's mean profile, selected
independently of it. Two variants: \textbf{none} (the default --- exactly zero perturbation, not
merely a small one) and \textbf{Dryden}, the MIL-F-8785C/MIL-HDBK-1797 horizontal turbulence
spectrum. It is a tier-1 feature only: the measured tiers carry whatever variability their own
profiles contain and nothing synthetic is added on top of them, since a sounding is already a
measurement of the real field rather than a model of its mean.

Scale length and intensity come from the spec's low-altitude formulas, which are written in feet
--- so $h_{ft}=h/0.3048$ is formed once, with $h$ the AGL altitude clamped at zero, and the band
logic stays in feet (a foot-valued threshold has an exact decimal value in feet and not in metres):
\begin{equation} \label{eq:dryden_scales}
    L_u = L_v = \frac{h_{ft}}{\left(0.177+0.000823\,h_{ft}\right)^{1.2}}\ \mathrm{[ft]}, \qquad
    \sigma_u = \sigma_v = \frac{0.1\,W_{20}}{\left(0.177+0.000823\,h_{ft}\right)^{0.4}},
\end{equation}
with $W_{20}$ the same mean profile \autoref{eq:wind_shear} evaluated at $20\,\mathrm{ft}$ --- one
expression for the mean wind, so a run cannot have two of them. $\sigma$ needs no unit conversion
(the formula is linear in $W_{20}$ and the bracket is dimensionless); $L$ does, on the way out.
Three bands:
\begin{center}
\begin{tabular}{l l l l}
Band & Condition & $L_u=L_v$ & $\sigma_u=\sigma_v$ \\
low        & $h_{ft}\leq1000$      & \autoref{eq:dryden_scales} & \autoref{eq:dryden_scales} \\
transition & $1000<h_{ft}<2000$    & linear, $1000\to1750\,\mathrm{ft}$ & linear, $0.1W_{20}\to\sigma_{high}$ \\
isotropic  & $h_{ft}\geq2000$      & $1750\,\mathrm{ft}$ & $\sigma_{high}$ (an input) \\
\end{tabular}
\end{center}
The low formula is exact at its own upper edge (the bracket is exactly $1$ at
$h_{ft}=1000$), so the transition band is continuous at both ends by construction. That band is an
implementor's fill, not a spec formula: the spec defines the low-altitude regime below
$1000\,\mathrm{ft}$ and the isotropic regime above $2000\,\mathrm{ft}$ and leaves the gap
undefined, and what is interpolated across it is the two \emph{parameters} $(L,\sigma)$ --- the gap
is a gap in the parameters, and the amplitudes below go as $\sqrt{\Phi}$, so interpolating them
instead would not even give the same answer. $L$ and $\sigma$ are continuous with a slope kink at
each boundary, the same kind of kink the ISA's piecewise lapse rates already put in $T(h)$ and
hence in the aerodynamic force; it is not smoothed, since no source specifies a smoothing.

The one-sided temporal power spectral densities are
\begin{gather}
    \Phi_u(\omega) = \sigma_u^2\,\frac{2L_u}{\pi v_{ref}}\,\frac{1}{1+\left(L_u\omega/v_{ref}\right)^2}
        \label{eq:dryden_psd_u} \\
    \Phi_v(\omega) = \sigma_v^2\,\frac{L_v}{\pi v_{ref}}\,
        \frac{1+3\left(L_v\omega/v_{ref}\right)^2}{\left[1+\left(L_v\omega/v_{ref}\right)^2\right]^2}
        \label{eq:dryden_psd_v}
\end{gather}
both normalised so that $\int_0^\infty\Phi\,\mathrm{d}\omega=\sigma^2$ exactly. They are realised
as a frozen random-phase sum over a log-spaced frequency grid $\{\omega_i\}$ with bin widths
$\{\Delta\omega_i\}$ (periods from $600\,\mathrm{s}$, longer than any flight, down to
$0.2\,\mathrm{s}$, which the break frequency $v_{ref}/L_u$ demands just off the pad):
\begin{equation} \label{eq:dryden_synthesis}
    k_i = \Phi(\omega_i)\Delta\omega_i, \qquad
    w_i = \sqrt{\frac{2k_i}{\sum_jk_j}}, \qquad
    u(t) = \sigma_u\sum_i w_i\sin(\omega_it+\varphi_i),
\end{equation}
and likewise for $v(t)$ with its own independent phases. The textbook amplitude is
$\sqrt{2\Phi_i\Delta\omega_i}$; the normalisation by $\sum_jk_j$ rescales it so that $\sum_iw_i^2=2$
and hence $\mathrm{Var}[u]=\sigma_u^2$ \emph{exactly}, for any band and any grid, instead of the
$81$--$98\%$ of it a finite band captures. That keeps $\sigma$ --- the quantity the spec actually
specifies --- as the physical knob, and leaves the band endpoints as pure numerics. The
$\{\omega_i\}$ are a deterministic grid and the phases $\{\varphi_i\}$ are drawn once per seed and
frozen, for the reason the next paragraph gives.

The two components are, by construction, \emph{along} the mean wind ($u$) and horizontally
\emph{across} it ($v$), so they compose with the flight-plane projection above rather than being
added to $W$ as scalars. With $\hat{e}_u$ the unit vector along the mean flow, $\hat{e}_v$ it
rotated by $+90^\circ$ in the horizontal plane, and $\hat{y}_0$ the cross-plane axis (to the right of downrange),
\begin{equation} \label{eq:dryden_composition}
    \hat{e}_u = -\cos(\psi-\chi)\,\hat{x}_0 - \sin(\psi-\chi)\,\hat{y}_0, \qquad
    \hat{e}_v = \sin(\psi-\chi)\,\hat{x}_0 - \cos(\psi-\chi)\,\hat{y}_0,
\end{equation}
\begin{equation} \label{eq:dryden_wind}
    (w_x,w_y) = (W+u)\,\hat{e}_u + v\,\hat{e}_v,
\end{equation}
which reduces to the mean-only projection above exactly when $u=v=0$. Note what this does
\emph{not} do: there is no non-negativity clamp on $W+u$. Perturbing a vector along $\hat{e}_u$
makes a negative $W+u$ the correct statement that the gust has reversed the local flow, where
clamping a perturbed \emph{speed} would silently flip the direction instead. And note what it
cannot do yet: the 2D closure discards $w_y$, so $v$ has no effect on any trajectory this model
integrates; it is carried so that a later 6DOF model picks it up unchanged.

Two deliberate limitations. First, $v_{ref}$ is a \emph{fixed input},
not the rocket's own relative airspeed $v_{rel}$: Dryden's temporal spectra
describe what a vehicle flying at a constant airspeed through a frozen turbulence field measures,
and this rocket's airspeed spans zero on the pad to supersonic at burnout. Driving the spectrum
with the instantaneous airspeed is incompatible with the determinism the adaptive solver needs ---
the airspeed $v_{rel}$ at time $t$ is part of what the solver is still solving for, and it re-evaluates the
right-hand side at arbitrary, repeated and rejected times, so a quantity that is not a function of
$(t,\vec{y})$ alone destroys its step-size control and its reproducibility. The damage is bounded
and worth stating precisely: after the renormalisation in \autoref{eq:dryden_synthesis} the
prefactors cancel, so $v_{ref}$ and $L$ enter only through the ratio $L/v_{ref}$ --- a time, the
eddy-crossing period. $v_{ref}$ therefore has \emph{no} effect on the delivered intensity, only on
where the spectrum's knee sits in frequency, i.e.\ whether the rocket meets 30-second swells or
one-second buffets. The practical rule: set it to the airspeed over the part of the flight whose
turbulence matters to the question being asked. Lifting this needs the sampled evaluation mode and
a filter formulation to run on it. Second, there is no vertical component ($w$), as for the mean
profile. At $h=0$ the low formula gives $L=0$ and the spectrum collapses to zero energy
while $\sigma$ stays finite; the perturbation is taken as exactly zero there, which is the right
limit --- at the surface there is no eddy left to resolve.

\subsubsection{Vertical wind}\label{sec:vertical_wind}
All tiers above return horizontal wind only ($w_z=0$); updrafts and downdrafts are not modelled.
The structure already admits them: $\vec{v}_w=w_x\hat{x}_0+w_z\hat{z}_0$ in $\mathcal{F}_0$ with
$\hat{z}_0$ pointing down, so an updraft of speed $w_{up}$ is $w_z=-w_{up}$ and changes
$\vec{v}_{rel}=\vec{v}_r-\vec{v}_w$ (hence $v_{rel}$ and $\alpha$) directly. Its size matters
most where $v_{rel}$ is small --- near the pad and around apogee, and under parachute --- since
a vertical wind $w_{up}$ rotates the relative wind by about $w_{up}/v_{rel}$. Two standard families exist in the literature for deriving a
vertical component: (1) terrain-slope / mass-continuity orographic lift, $w\approx\vec{v}_h\cdot
\nabla h_{terrain}$, which needs a terrain elevation model this project does not have; (2)
thermal/convective updraft parameterisations (as used in the soaring/glider literature), which need
surface heat-flux or insolation data outside the radiosonde soundings used here. Neither fits the
current data pipeline (station soundings + a simple analytic profile) without adding a new data
source. This is deliberately left unimplemented for now, pending a look at which of the two
families (or a combination) is actually worth the added data dependency; revisit once that
research is done.

\subsection{Relative wind and angle of attack}\label{sec:relwind}
Let $\vec{v}_r$ be the rocket's velocity relative to $\mathcal{F}_0$, and $\vec{v}_w$ the wind
velocity relative to $\mathcal{F}_0$ (a model output). The rocket's velocity
relative to the surrounding air is
\begin{equation}
    \vec{v}_{rel} = \vec{v}_r - \vec{v}_w = v_{rel,x}\hat{x}_r + v_{rel,z}\hat{z}_r,
\end{equation}
with $v_{rel,x}$ the axial (along-nose) component and $v_{rel,z}$ the normal component, both
expressed in $\mathcal{F}_r$. The relative airspeed and angle of attack are
\begin{equation}
    v_{rel} = \sqrt{v_{rel,x}^2+v_{rel,z}^2}, \qquad \alpha = \operatorname{atan2}(v_{rel,z}, v_{rel,x}),
\end{equation}
so that $v_{rel,x}=v_{rel}\cos\alpha$, $v_{rel,z}=v_{rel}\sin\alpha$. At $\alpha=0$ (nose-first into the relative wind)
$\vec{v}_{rel}=v_{rel}\hat{x}_r$.

Define the unit vector along the relative wind and its in-plane normal (a $+90^\circ$ rotation of
$\hat{u}$ within $\mathcal{F}_r$):
\begin{equation}
    \hat{u} = \frac{\vec{v}_{rel}}{v_{rel}} =\cos\alpha\,\hat{x}_r + \sin\alpha\,\hat{z}_r, \qquad
    \hat{n} = -\sin\alpha\,\hat{x}_r + \cos\alpha\,\hat{z}_r.
\end{equation}
At $\alpha=0$: $\hat{u}=\hat{x}_r$, $\hat{n}=\hat{z}_r$. $\hat{u},\hat{n}$ feed the dynamic pressure
$q$ and, through $\alpha$, the coefficient lookup; the aerodynamic force itself is resolved along
the fixed body axes $\hat{x}_r,\hat{z}_r$, \emph{not} along $\hat{u},\hat{n}$.

\subsection{Non-inertial frame corrections (Earth's rotation)}\label{sec:noninertial}
The WGS84 gravity model above already accounts for the Earth's rotation in the magnitude of $g$
(the static, equilibrium-shape effect, folded into $\gamma_0(\phi)$). A rotating Earth also
introduces two velocity/position-dependent fictitious forces into the translational equations of
motion, once $\mathcal{F}_0$ is reinterpreted as Earth-fixed rather than inertial:
\begin{equation}
    \vec{a}_{Coriolis} = -2\vec{\Omega}\times\vec{v}_r, \qquad
    \vec{a}_{centrifugal} = -\vec{\Omega}\times(\vec{\Omega}\times\vec{r}),
\end{equation}
with $\vec{\Omega}$ the Earth's angular velocity ($|\vec{\Omega}|=7.2921\times10^{-5}\,
\mathrm{rad/s}$) and $\vec{r},\vec{v}_r$ the rocket's position/velocity relative to the Earth-fixed
frame. The centrifugal term is a small, latitude-dependent correction, already absorbed into the
WGS84 $g(\phi,h)$ above when the launch site is the reference point; the Coriolis term is
genuinely new here and depends on the flight's own velocity and heading, not just the launch site.

\paragraph{Ground-frame basis.} Work in the local East--North--Up triad $(\hat{E},\hat{N},\hat{U})$
at the launch site $(\phi_0,\lambda_0)$, with $\phi_0$ the site's geodetic latitude (the same
quantity written $\phi$ in the gravity model). The ground frame $\mathcal{F}_0$ is fixed by the
launch heading $\chi$ (degrees clockwise from North, the direction of $+\hat{x}_0$)
and by $\hat{z}_0$ pointing \emph{down}:
\begin{equation}
    \hat{x}_0 = \sin\chi\,\hat{E} + \cos\chi\,\hat{N}, \qquad
    \hat{y}_0 = \cos\chi\,\hat{E} - \sin\chi\,\hat{N}, \qquad
    \hat{z}_0 = -\hat{U}.
\end{equation}
This whole subsection is written in the \emph{fixed} ground frame $\mathcal{F}_0$, not in the
(possibly rotating) gravity direction $\hat{z}_g$: the $\vec{\Omega}$ decomposition below is
evaluated once, at the launch site, and does \emph{not} follow a rotating $\hat{z}_g(\varphi)$ ---
see the closing paragraph for what that costs. The triad is
right-handed, $\hat{x}_0\times\hat{y}_0=\hat{z}_0$, which can be checked directly from the three
lines above ($(\sin\chi,\cos\chi,0)\times(\cos\chi,-\sin\chi,0)=(0,0,-1)$ in ENU components).
$\hat{y}_0$ points to the \emph{right} of the downrange direction.

\paragraph{Components of $\vec{\Omega}$.} The Earth's angular velocity points along the rotation
axis, i.e.\ toward the celestial north pole. At latitude $\phi_0$ that direction lies in the local
meridian plane (no East component), tilted from the local vertical by the co-latitude:
\begin{equation}
    \vec{\Omega} = \Omega\left(\cos\phi_0\,\hat{N} + \sin\phi_0\,\hat{U}\right), \qquad
    \Omega = 7.2921\times10^{-5}\,\mathrm{rad/s}.
\end{equation}
Projecting onto the ground-frame triad by taking dot products with the three basis vectors above:
\begin{equation}
    \Omega_x = \vec{\Omega}\cdot\hat{x}_0 = \Omega\cos\phi_0\cos\chi, \qquad
    \Omega_y = \vec{\Omega}\cdot\hat{y}_0 = -\Omega\cos\phi_0\sin\chi, \qquad
    \Omega_z = \vec{\Omega}\cdot\hat{z}_0 = -\Omega\sin\phi_0,
\end{equation}
so that
\begin{equation} \label{eq:omega_components}
    \vec{\Omega} = \Omega\left(\cos\phi_0\cos\chi\,\hat{x}_0
        - \cos\phi_0\sin\chi\,\hat{y}_0 - \sin\phi_0\,\hat{z}_0\right),
\end{equation}
with $\Omega_x^2+\Omega_y^2+\Omega_z^2=\Omega^2$ as required. Three sanity checks: at the North
pole ($\phi_0=90^\circ$) it reduces to $-\Omega\hat{z}_0$, i.e.\ straight up, for every $\chi$; on
the equator launching due North ($\phi_0=0,\chi=0$) to $+\Omega\hat{x}_0$, i.e.\ horizontal and
downrange; on the equator launching due East ($\phi_0=0,\chi=90^\circ$) to $-\Omega\hat{y}_0$,
which is horizontal and to the left of downrange --- North again, as it must be. The sign of
$\Omega_z$ is purely a bookkeeping consequence of this project's $z$-down convention: the rotation
vector points \emph{up} out of the northern hemisphere, hence negative along $\hat{z}_0$.

\paragraph{Which component of $\vec{\Omega}$ actually matters in 2D.} With motion confined to the
flight plane, $\vec{v}_r=v_x\hat{x}_0+v_z\hat{z}_0$ (no $\hat{y}_0$ component). Expanding the cross
product in the right-handed triad,
\begin{equation}
    \vec{\Omega}\times\vec{v}_r =
        \Omega_yv_z\,\hat{x}_0
        + \left(\Omega_zv_x - \Omega_xv_z\right)\hat{y}_0
        - \Omega_yv_x\,\hat{z}_0,
\end{equation}
so
\begin{equation} \label{eq:coriolis_split}
    \vec{a}_{Coriolis} = -2\vec{\Omega}\times\vec{v}_r =
        \underbrace{-2\Omega_y\left(v_z\hat{x}_0 - v_x\hat{z}_0\right)}_{\text{in-plane}}
        \;\underbrace{-\,2\left(\Omega_zv_x - \Omega_xv_z\right)\hat{y}_0}_{\text{out-of-plane}}.
\end{equation}
This is worth reading carefully, because it inverts the intuitive move. The in-plane Coriolis
acceleration is produced \emph{entirely} by $\Omega_y$ --- the component of $\vec{\Omega}$
\emph{normal} to the flight plane --- while $\Omega_x$ and $\Omega_z$, the components lying
\emph{inside} the plane, produce a purely out-of-plane acceleration. ``Project $\vec{\Omega}$ onto
the flight plane and keep that part'' would therefore keep exactly the half that does nothing
in-plane and throw away the half that does everything.

\paragraph{Simplification (2D closure).} Only the in-plane bracket of \autoref{eq:coriolis_split}
is retained; the $\hat{y}_0$ term is dropped. This is \emph{not} a smallness argument --- the
discarded lateral acceleration is of the same order as the retained one, and for some geometries
larger. It is dropped because the model has no lateral degree of freedom to receive it: the
out-of-plane Coriolis deflection (the familiar right-hand deflection in the northern hemisphere) is
physically real, and a 2D planar model simply cannot represent it. This is the same closure already
applied to the crosswind component $w_y$, and it is lifted by the same future work (a full 6DOF
model).

Two consequences follow immediately and should be stated so results are not over-read: $\Omega_y=
-\Omega\cos\phi_0\sin\chi$ vanishes for a launch due North or due South ($\chi=0^\circ,180^\circ$)
and at the poles, so the 2D model predicts \emph{zero} Coriolis effect there while the true
Coriolis effect is maximal and entirely lateral; and $|\Omega_y|$ is largest for an East/West
launch near the equator, where the effect the 2D model \emph{can} see is also the largest it ever
is.

\paragraph{In-plane Coriolis term.} Substituting $\Omega_y$ into the retained bracket,
\begin{equation} \label{eq:coriolis_2d}
    \vec{a}_{Coriolis}^{\,2D} = 2\Omega\cos\phi_0\sin\chi\left(v_z\,\hat{x}_0 - v_x\,\hat{z}_0\right),
\end{equation}
i.e.\ componentwise ($v_x$ the downrange component of $\vec{v}_r$,
$v_z$ the component \emph{downward}, so an ascending rocket has $v_z<0$):
\begin{equation}
    a_{Coriolis,x} = +2\Omega\cos\phi_0\sin\chi\;v_z, \qquad
    a_{Coriolis,z} = -2\Omega\cos\phi_0\sin\chi\;v_x.
\end{equation}
The two classical checks both come out right. A rocket rising vertically at the equator with an
Eastward $+\hat{x}_0$ ($\phi_0=0$, $\chi=90^\circ$, $v_z<0$) gets $a_x<0$: it drifts
\emph{West}, the standard counterpart of the eastward deflection of a falling body. A body moving
East ($v_x>0$) gets $a_z<0$, i.e.\ \emph{upward} --- the E\"otv\"os effect, eastward motion
reducing apparent weight.

\paragraph{Residual centrifugal term.} Write the rocket's geocentric position as
$\vec{r}=\vec{R}_0+\vec{s}$, with $\vec{R}_0$ from the Earth's centre to the launch site and
$\vec{s}=x\hat{x}_0+z\hat{z}_0$ the simulation's own position state. Since
$-\vec{\Omega}\times(\vec{\Omega}\times\vec{a})=\Omega^2\vec{a}-(\vec{\Omega}\cdot\vec{a})
\vec{\Omega}$ is linear in $\vec{a}$, the centrifugal acceleration splits cleanly into a constant
part, already carried by the WGS84 Somigliana formula, and a residual:
\begin{equation}
    \vec{a}_{centrifugal} = \underbrace{\Omega^2\vec{R}_0-(\vec{\Omega}\cdot\vec{R}_0)\vec{\Omega}}
        _{\text{constant, already in }\gamma_0(\phi)}
        + \underbrace{\Omega^2\vec{s}-(\vec{\Omega}\cdot\vec{s})\vec{\Omega}}_{\text{residual}}.
\end{equation}
This residual is bounded by $\Omega^2|\vec{s}|=5.3\times10^{-9}\,|\vec{s}|$, i.e.\
$5\times10^{-5}\,\mathrm{m/s^2}$ even at $|\vec{s}|=10\,\mathrm{km}$ --- roughly three orders of
magnitude below the Coriolis term, and its \emph{vertical} part is in any case already folded into
the WGS84 free-air correction through the constant $m$. Only the Coriolis term is implemented; the
residual centrifugal term is left unimplemented.

\paragraph{Moment.} Coriolis and centrifugal accelerations are body forces distributed in exact
proportion to mass, so their resultant acts at the centre of gravity $g$ and contributes
\emph{zero} moment about $g$. The one surviving coupling comes from the body's own pitch rate
($\vec{v}_i=\vec{v}_g+\vec{\omega}_r\times\vec{r}_i$ inside the rigid body), a gyroscopic torque of
order $2I\Omega\omega_r$, i.e.\ an angular acceleration of order $2\Omega\omega_r\approx
1.5\times10^{-4}\,\mathrm{rad/s^2}$ at $\omega_r=1\,\mathrm{rad/s}$. This is negligible against the
aerodynamic pitching moment and is dropped; the rotational equation is unchanged.

\paragraph{Result.} The translational equation gains one term:
\begin{equation}
    m\vec{a}_r = \vec{W} +\vec{T} + \vec{A}_f + m\,\vec{a}_{Coriolis}^{\,2D},
\end{equation}
with $\vec{a}_{Coriolis}^{\,2D}$ from \autoref{eq:coriolis_2d}, evaluated directly in
$\mathcal{F}_0$ (no body-axis rotation needed --- it is already a ground-frame expression, like
$\vec{W}$). No new parameter is introduced beyond the launch-site latitude $\phi_0$ and
heading $\chi$.

\paragraph{Magnitudes.} The retained term scales as
$2\Omega\cos\phi_0\sin\chi\,|\vec{v}_r|$: at $|\vec{v}_r|=300\,\mathrm{m/s}$ it is at most
$4.4\times10^{-2}\,\mathrm{m/s^2}$ (equatorial East launch) and about
$2.8\times10^{-2}\,\mathrm{m/s^2}$ at $\phi_0=50^\circ$ --- some $0.3\%$ of $g$, accumulating to
order $10\,\mathrm{m}$ of displacement over a $30\,\mathrm{s}$ powered ascent. It is roughly the
same order as the effective $g\,x/R_e\approx1.5\times10^{-2}\,\mathrm{m/s^2}$ (at
$x=10\,\mathrm{km}$) that gravity acquires by pointing along the rotated local vertical; the two
remain \emph{independent} and must not be multiplied together.

\paragraph{One residual inconsistency.} The decomposition
\autoref{eq:omega_components} is evaluated once, at the launch site's latitude $\phi_0$ and
heading $\chi$, and is \emph{not} re-evaluated as the rocket travels downrange --- so
$\vec{\Omega}$ is resolved in the fixed ground frame $\mathcal{F}_0$ while gravity, if the
radial/elliptical direction model is enabled, is resolved along the rotated $\hat{x}_g,\hat{z}_g$.
Deliberately so: building a single consistently rotating frame would replace both corrections with
a curvilinear formulation the 2D state cannot carry (it would need the metric terms of motion over
a sphere, not just a rotated basis). The cost of leaving them independent is bounded and small.
Along the ground track the latitude changes by $x\cos\chi/R_e$ and the heading by the meridian
convergence, so $\Omega_y$ drifts by $O(\Omega\,x/R_e)$, i.e.\ a relative $3\times10^{-6}$ of the
retained Coriolis acceleration at $x=10\,\mathrm{km}$ and $3\times10^{-4}$ at $100\,\mathrm{km}$
--- five orders of magnitude smaller than the gravity-direction term the same $x/R_e$ produces,
because the Coriolis term is itself $10^{-3}$ of gravity to begin with.

\subsection{States, initial conditions and integration}\label{sec:states}
This subsection fixes what is integrated, from where, and how. The rocket has two groups of
kinematic states, both resolved in the fixed ground frame $\mathcal{F}_0$:
\begin{itemize}
    \item \emph{translation}: position $\vec{x}_r=x\,\hat{x}_0+z\,\hat{z}_0$, velocity
    $\vec{v}_r=v_x\,\hat{x}_0+v_z\,\hat{z}_0$ and acceleration $\vec{a}_r$;
    \item \emph{orientation}: pitch angle $\theta_r$, pitch rate $\omega_r$ and angular acceleration
    $\varepsilon_r$.
\end{itemize}
Each level is the time derivative of the one below it,
\begin{equation} \label{eq:state_kinematics}
    \frac{\mathrm{d}\vec{x}_r}{\mathrm{d}t}=\vec{v}_r, \quad
    \frac{\mathrm{d}\vec{v}_r}{\mathrm{d}t}=\vec{a}_r, \qquad
    \frac{\mathrm{d}\theta_r}{\mathrm{d}t}=\omega_r, \quad
    \frac{\mathrm{d}\omega_r}{\mathrm{d}t}=\varepsilon_r, \qquad
    \frac{\mathrm{d}m}{\mathrm{d}t}=-\frac{T_0(t)}{\mathrm{ISP}\,g_0},
\end{equation}
and the accelerations $\vec{a}_r$, $\varepsilon_r$ are not independent states: they are the outputs
of \autoref{eq:eom_trans} and \autoref{eq:eom_rot}, evaluated from the other states and time at
every call of the right-hand side. The last relation is the propellant mass flow. The integrated
state vector is therefore of dimension seven,
\begin{equation} \label{eq:state_vector}
    \vec{y} = \left(x,\ z,\ v_x,\ v_z,\ \theta_r,\ \omega_r,\ m\right),
\end{equation}
with $z$ and $v_z$ positive \emph{down}. The mass $m$ is only read when the variable-mass model is
switched on; otherwise $m$ and $I$ are fixed and $\mathrm{d}m/\mathrm{d}t=0$. The right-hand side
has three modes, selected from the state in this order: on the rail, under canopy, otherwise free
flight.
\todo{The code evaluates the mass flow with the thrust magnitude $T$ (including the nozzle-exit
pressure term), not the reference thrust $T_0$; the two differ only when the exit area is non-zero.}

\paragraph{Initial conditions.} The default start is at rest on the pad, at ignition, with the nose
along the rail:
\begin{equation} \label{eq:initial_conditions}
    \vec{x}_r(0)=\vec{0}, \quad \vec{v}_r(0)=\vec{0}, \quad
    \theta_r(0)=\theta_{rail}, \quad \omega_r(0)=0, \quad m(0)=m_{launch},
\end{equation}
with $\theta_{rail}$ the rail elevation (default $90^\circ$) and $m_{launch}=m_r+m_0$ the liftoff
mass. Time starts at $t=0$, the ignition instant of the thrust curve. The launch site supplies
\begin{itemize}
    \item latitude $\phi_0$ and heading $\chi$ --- Earth-rotation correction and ellipsoidal gravity;
    \item elevation $h_0$ --- atmosphere and gravity only; it never enters $z$;
    \item rail elevation and length --- the launch rail.
\end{itemize}
A scenario may override the whole state; nothing checks that $\theta_r(0)=\theta_{rail}$. The
models carry no hidden dynamical state, so they are pure functions of $(t,\vec{y})$.

\paragraph{Integrators.} Two integrators are available behind one interface, chosen by a setting:
\begin{itemize}
    \item \textbf{Adaptive Runge--Kutta 4(5)} (the default): the \textsc{scipy} \texttt{RK45}
    scheme with relative tolerance $10^{-8}$ and absolute tolerance $10^{-10}$ on every component of
    $\vec{y}$, dense output on, and a time limit of $300\,\mathrm{s}$. Events (below) are located
    to solver tolerance by root finding on the dense output.
    \item \textbf{Fixed-step classical Runge--Kutta 4}, a hand-written scheme with step $h$
    (default $1\,\mathrm{ms}$):
    \begin{equation} \label{eq:rk4}
        \vec{y}_{k+1}=\vec{y}_k+\frac{h}{6}\left(\vec{k}_1+2\vec{k}_2+2\vec{k}_3+\vec{k}_4\right),
    \end{equation}
    with $\vec{k}_1=\vec{f}(t_k,\vec{y}_k)$, $\vec{k}_{2,3}=\vec{f}\!\left(t_k+\tfrac{h}{2},
    \vec{y}_k+\tfrac{h}{2}\vec{k}_{1,2}\right)$, $\vec{k}_4=\vec{f}(t_k+h,
    \vec{y}_k+h\,\vec{k}_3)$ and $\vec{f}$ the right-hand side. It has no step control and no
    event location: events are detected afterwards from sign changes between consecutive steps and
    placed by linear interpolation inside the step, so they are accurate to about $h$. It is
    the cross-check of the adaptive path, and the only one that supports the sampled-data mode.
\end{itemize}
Both integrate the same right-hand side, so the two differ only in the numerical error.
\todo{Document the cost/accuracy trade-off between the two integrators with a convergence
study (error against $h$ and against tolerance).}

\paragraph{Events.} Five events are watched; only the first ends the run. Each is a sign change of
a scalar function of $(t,\vec{y})$:
\begin{itemize}
    \item \textbf{ground impact} (terminal): the ground depth crosses zero going down, $z$ for the
    flat ground, the curved-ground depth otherwise. The run stops here (or at the time limit).
    Ignored during the first $0.1\,\mathrm{s}$, since $z=0$ at the start;
    \item \textbf{apogee}: $v_z$ crosses from negative to positive in the launch frame, again
    ignored for the first $0.1\,\mathrm{s}$;
    \item \textbf{burnout}: the thrust magnitude falls through a small threshold ($10^{-9}$\,N)
    from above, so a motor that never produces thrust never triggers it;
    \item \textbf{rail exit}: $s_{rail}$ crosses $L_{rail}$ upwards;
    never active when $L_{rail}=0$;
    \item \textbf{parachute deployment}: the deployment predicate turns true; never active without a recovery model.
\end{itemize}
Apart from ground impact the events are informational: they only report times, and none of them
switches the dynamics. The rail, canopy and air-brake modes are state predicates evaluated inside
the right-hand side (the right-hand side cannot read an event that has not been finalised yet),
which is why the same condition appears twice: as a predicate that acts and as an event that
reports.
\todo{The apogee event is the launch-frame $v_z$ crossing, not the maximum of the true altitude
above a curved ground; redefining apogee is a separate modelling decision.}

\subsection{Launch rail}\label{sec:rail}
While the rocket is still on the launch rail it can slide along the rail and do nothing else: the
rail buttons mechanically react every force component transverse to the rail and prevent rotation.
This is a \emph{kinematic constraint}, not a new force: the already-assembled acceleration is
projected onto the rail, which is the same as the buttons reacting exactly the transverse part of
the force, and the rotation is removed together with it.

The rail direction is fixed for the whole flight, and is the direction the nose initially points.
With $\theta_{rail}$ the rail's elevation above the horizontal, an as-built input to be set
consistently with the initial $\theta_r$,
\begin{equation} \label{eq:rail_direction}
    \hat{e}_{rail} = \cos\theta_{rail}\,\hat{x}_0 - \sin\theta_{rail}\,\hat{z}_0,
\end{equation}
that is, $\hat{x}_r$ evaluated at $\theta_r=\theta_{rail}$ and then frozen. The distance travelled along the rail, from a launch point at the origin, is the
projection of the position state onto it,
\begin{equation} \label{eq:rail_distance}
    s_{rail}(\vec{x}_r) = \vec{x}_r\cdot\hat{e}_{rail}
        = x\cos\theta_{rail}-z\sin\theta_{rail},
\end{equation}
which is $0$ at $t=0$ and grows as the rocket slides. The rocket is on the rail while
$s_{rail}<L_{rail}$ (a state predicate, re-evaluated at every call of the right-hand side), and
there the equations of motion are replaced by
\begin{gather}
    \vec{a}_r = \left[\left(\frac{\sum\vec{F}}{m}
        + \vec{a}_{Coriolis}^{\,2D}\right)\cdot\hat{e}_{rail}\right]\hat{e}_{rail}
        \label{eq:rail_constraint} \\
    \omega_r = 0, \qquad \varepsilon_r = 0 \label{eq:rail_attitude}
\end{gather}
with $\sum\vec{F}$ exactly the sum of \autoref{eq:eom_trans} and $\vec{a}_{Coriolis}^{\,2D}$ the
term of \autoref{eq:coriolis_2d} if it is enabled. The component of $\vec{a}_r$ normal to the rail
is therefore identically zero, and with $\vec{x}_r(0)$ on the rail line and $\vec{v}_r(0)=\vec{0}$
the position and velocity stay on that line for the whole phase, so the constraint needs no
correction step. In the integrator, \autoref{eq:rail_attitude} reads $\mathrm{d}\theta_r/\mathrm{d}t=0$,
$\mathrm{d}\omega_r/\mathrm{d}t=0$, i.e.\ $\theta_r$ and $\omega_r$ keep their initial values.
The constraint reacts the \emph{same} total acceleration every other model produces. The
pitching moment is reacted by the rail rather than allowed to act.

\paragraph{Exit and reverse motion.} $L_{rail}=0$ is the default and the off switch. Past rail
exit ($s_{rail}\geq L_{rail}$) the free-flight equations apply again, with no transition term.
Nothing in the model holds the rocket on the pad before thrust exceeds weight: the rail
projection of \autoref{eq:rail_constraint} is applied as it stands, whatever its sign.
\todo{No pad contact: if the thrust along the rail is below the weight component, \autoref{eq:rail_constraint}
gives a negative $\vec{a}_r\cdot\hat{e}_{rail}$ and the rocket slides \emph{backwards} down the
rail ($s_{rail}<0$) until thrust catches up. A real pad reacts weight until lift-off; the fix is to
clamp $\vec{a}_r\cdot\hat{e}_{rail}\geq0$ while $s_{rail}\leq0$ and $\vec{v}_r\cdot\hat{e}_{rail}\leq0$,
which would change existing results for low thrust-to-weight starts.}
The physically interesting output is the rail-exit speed $|\vec{v}_r|$, which decides whether the
wind can tip the airframe before it has the airspeed to weathercock.

Simplified: no rail friction, no rail flex, and no tip-off --- the transverse and rotational
degrees of freedom are restored instantaneously at $s_{rail}=L_{rail}$.

\subsection{Recovery}\label{sec:recovery}
The second mechanically constrained phase, and --- usefully --- the same mechanism as the first:
once a recovery canopy is deployed the attitude state is frozen exactly as on the rail,
$\omega_r=0$, $\varepsilon_r=0$, and only the translational equation changes. Deployment is
instantaneous and total, and the canopy is modelled by its drag alone. During the descent the
translation obeys \autoref{eq:eom_trans} with one extra force,
\begin{equation} \label{eq:chute_eom}
    m\vec{a}_r = \vec{W} + \vec{T} + \vec{A}_f + \vec{F}_{ch},
\end{equation}
with $\vec{F}_{ch}$ the canopy drag of \autoref{eq:chute_force}. The orientation is held at
the pitch angle $\theta_r$ it had at deployment, and the pitching moment is not integrated.
\todo{In the code only $\mathrm{d}\theta_r/\mathrm{d}t$ and $\mathrm{d}\omega_r/\mathrm{d}t$ are set
to zero, so the stored $\omega_r$ keeps its non-zero value at deployment and the link
$\omega_r=\mathrm{d}\theta_r/\mathrm{d}t$ is broken under canopy.
The reported pitch rate is then meaningless after deployment; consider resetting it to zero.}
Freezing is chosen over an artificial damping term, which would need a time constant no source
fixes.

\paragraph{Deployment.} A function of the state being integrated, not of a detected event ---
either a timer, $t\geq t_{dep}$, or apogee, taken as $v_z>0$ (the rocket is descending, since
$z$ is positive \emph{down}). The second one deserves its caveat stated: ``has apogee happened''
is reconstructed from the state, because the right-hand side is evaluated before any event of the
solve is finalised and cannot consult one, so a pre-apogee transient in $v_z$ --- which an
unstable trajectory can produce --- would arm deployment early. Not guarded against; recorded.
These two triggers are the only ones; the predicate is not latched, but both are monotone on a
nominal flight. While the rocket is still on the rail the rail mode takes precedence.
\todo{No altitude (or velocity/dynamic-pressure) deployment trigger exists; a main-chute-at-altitude
trigger and a drogue-plus-main sequence are not modelled.}

\paragraph{Force.} The canopy's drag is computed directly in $\mathcal{F}_0$, like weight and
unlike the body-axis aerodynamic force:
\begin{equation} \label{eq:chute_force}
    \vec{F}_{ch} = -\tfrac12\rho\,C_{D,ch}S_{ch}\,v_{rel}\,\vec{v}_{rel},
\end{equation}
with $\vec{v}_{rel}=\vec{v}_r-\vec{v}_w$ and $v_{rel}=|\vec{v}_{rel}|$,
$S_{ch}$ the canopy reference area and $C_{D,ch}$ its drag coefficient; the model is the drag area
$C_{D,ch}S_{ch}$ of the canopy, acting against the relative wind. Three properties of this equation are
the whole content of the model, and each is a decision rather than a convenience:
\begin{itemize}
    \item It is \emph{natively inertial}, and it is \emph{not} routed through the body
    axial/normal decomposition. A round canopy is axisymmetric and its drag
    opposes the relative-wind \emph{vector}, not the body's nose direction --- which, under canopy,
    is frozen at whatever attitude deployment happened to leave it. Routing it through
    $\hat{x}_r$ would make the descent rate depend on that attitude, which is not how a parachute
    behaves. $C_{D,ch}$ is therefore a drag coefficient on $\hat{v}_{rel}$ and specifically
    \emph{not} the axial $C_x$: it belongs to a different force, in a
    different frame, with a different reference area ($S_{ch}$, not $S$).
    \item It is \emph{added to}, not substituted for, the airframe's own forces. Weight, any
    residual thrust and the body-axis aerodynamic force all continue to act; $\vec{F}_{ch}$ is one
    more term in that sum, and the frozen $\theta_r$ is what they are evaluated at.
    \item It is written as $v_{rel}\,\vec{v}_{rel}$ rather than the algebraically identical
    $\tfrac12\rho v_{rel}^2C_{D,ch}S_{ch}\,\hat{v}_{rel}$, so that it has no $0/0$ direction at $v_{rel}=0$ --- the
    same device as \autoref{eq:pitch_damping_moment}.
\end{itemize}
The quantity worth checking a run against is the terminal descent speed, where $\vec{F}_{ch}$
balances weight:
\begin{equation} \label{eq:chute_terminal}
    v_{term} = \sqrt{\frac{2mg}{\rho\,C_{D,ch}S_{ch}}},
\end{equation}
approached rather than attained, since $\rho$ varies over the descent.

$S_{ch}=0$ is the default, and so is ``no recovery model'' --- either alone leaves this subsection
inert and reproduces every result produced without it. What is deliberately simplified: inflation
is instantaneous and total, with no opening transient, no inflation shock and no reefing stage; and
there is no pendulum swing of the body under the risers, which the frozen attitude replaces
outright.
\todo{Neglected and not yet in the model: the opening transient (drag area
growing from zero over an inflation time), the opening-shock load on the airframe and shock cord, a
reefed or two-stage (drogue and main) sequence, the pendulum and coning motion of the body under the
risers, and the added-mass effect of the air carried with the canopy.}

\subsection{Air brake}\label{sec:airbrake}
A switchable extra drag area, rigidly mounted to the body, used to modulate drag during the
unpowered coast. It stands alone, as an area and a drag coefficient of its own, rather than being
derived from an airframe.

It enters as one additive axial term, so that the axial force $qSC_x$ becomes $qSC_x+F_b$ with
\begin{equation} \label{eq:airbrake_force}
    F_b = q\,C_{D,b}\,S_b\,f_b, \qquad f_b\in[0,1],
\end{equation}
where $S_b$ is the exposed area at full deployment, $C_{D,b}$ its drag coefficient, and $f_b$ the
deployment fraction --- i.e.\ a linear fraction-to-area model, $S(f_b)=f_bS_b$, the simplest one
that needs no geometry beyond those two numbers. The term is routed along $\hat{x}_r$, since a
brake rigidly mounted to the body adds axial drag whatever the pitch. It contributes to the axial
force \emph{only}: nothing to $qSC_z$ or $qSc\,C_m$, and the arm $l_p-l_g$ does not move.
$S_b=0$ is the default, as is the ``never deploys'' law below.

Two orthogonal choices set $f_b$: a \emph{law}, which decides how far open the brake is, and a
\emph{trigger}, which decides from when it is allowed to open at all.

\subsubsection{Deployment law}\label{sec:airbrake_law}
Three interchangeable variants:
\begin{itemize}
    \item \textbf{None} (the default): $f_b\equiv0$. The brake never deploys, and every result
    produced without one is reproduced unchanged.
    \item \textbf{Proportional}: pure P control on \emph{predicted} apogee,
    \begin{equation} \label{eq:airbrake_p}
        f_b = \mathrm{clip}\!\left(k_p\left(\hat{h}_a-h_{tgt}\right),\,0,\,1\right), \qquad
        \hat{h}_a = h + \frac{v_z^2}{2g},
    \end{equation}
    with $h_{tgt}$ the commanded apogee and $k_p$ the gain. $\hat{h}_a$ is the standard estimator
    --- the altitude the rocket would reach if drag vanished at this instant ($v_z^2$ because
    the sign is squared away) --- and is deliberately an \emph{upper} bound: it ignores all further
    drag, including the brake's own. The error is positive when the rocket is predicted to
    overshoot, so more brake is commanded. The law is gated to the unpowered-ascent window,
    returning $0$ whenever $T(t)>0$ or $v_z\geq0$; that window is the law's own and is \emph{not}
    the trigger below.
    \item \textbf{Full}: $f_b=1$ from the trigger instant onward and $0$ before it --- a constant
    area, opened once. This is the variant a sizing study solves against: with $f_b$ a step, a
    target apogee fixes either the product $C_{D,b}S_b$ or the deployment time, one unknown against
    one equation, which the proportional law's closed loop does not offer.
\end{itemize}
Every law here is \emph{memoryless}: it is re-evaluated afresh from the current state at every
evaluation of the right-hand side, since the solver evaluates it at arbitrary, repeated and
rejected times. A law \emph{with} memory needs the sampled mode below; none is written.

\subsubsection{Deployment trigger}\label{sec:airbrake_trigger}
Which condition arms the brake, chosen independently of the law above. Three variants, each a pure
\emph{onset} latch --- once its condition is true it stays armed for the rest of the flight:
\begin{itemize}
    \item \textbf{After burnout} (the default): $T(t)\leq0$, the negation of the proportional law's
    own powered-flight gate, true from burnout (the falling zero-crossing of the thrust curve)
    onward.
    \item \textbf{After time}: $t\geq t_{brake}$, with $t_{brake}=0$ --- the default --- meaning
    deployed from $t=0$.
    \item \textbf{After apogee}: $v_z>0$, the same descending-state predicate as the canopy's
    apogee trigger, and carrying the same caveat for the same reason: it is a state predicate and
    never the apogee event, so a pre-apogee transient in $v_z$ would arm it early.
\end{itemize}
Only the \textbf{full} law actually reads the trigger; \textbf{none} is indifferent to it and
\textbf{proportional} uses its own hardcoded window instead, so selecting a trigger changes nothing
under those two. What no trigger expresses is \emph{retraction}: there is no condition under which
an armed brake closes again, and no deployment-rate limit between $f_b=0$ and $f_b=1$ --- a real
actuator has both, and both are future work.

\subsection{Sampled-data evaluation (zero-order hold)}\label{sec:sampled}
This subsection is an interchangeable \emph{mode of evaluating} the models, changing \emph{when} a
quantity is computed rather than \emph{what} it is.

By default every model is a pure function of $(t,\vec{y})$ (\autoref{eq:state_vector}),
re-evaluated at every call of the right-hand side. Nothing of the form ``computed once, then held
until the next tick'' --- a flight computer's measure--compute--hold cycle, an integral term, a
filter with internal state --- is expressible that way. This mode makes it expressible.

With a sample rate $f_s>0$ and period $T_s=1/f_s$, define the tick times $t_k=kT_s$. At each tick,
every registered \emph{channel} --- a named sampled quantity --- is evaluated once, at that tick's
own state, and held:
\begin{equation} \label{eq:zoh}
    c(t) = c_k \equiv c\!\left(t_k,\vec{y}(t_k);\,c_{k-1}\right), \qquad t\in[t_k,t_{k+1}).
\end{equation}
The argument worth noticing is the last one: a channel receives \emph{its own} previous held value,
$c_{k-1}$ (undefined at the first tick), rather than the whole previous state of the hold. That is
what makes a channel recursive, and therefore what makes an integral accumulator, a rate limiter's
previous command, or a discrete filter's internal state expressible at all --- each of them is
exactly a function of the form \autoref{eq:zoh}. One channel exists today, the deployment fraction
$f_b$; its present laws are memoryless and ignore $c_{k-1}$, so the
mode currently changes only \emph{how often} $f_b$ is computed, not what it computes.

Three requirements come with \autoref{eq:zoh}, and each is a consequence rather than a preference:
\begin{enumerate}
    \item \textbf{Fixed-step integration, with $T_s$ an exact integer multiple of the step.} An
    adaptive step cannot honour a hold boundary, so this mode requires the fixed-step scheme, and
    $T_s/h$ must be an integer so that every hold boundary falls exactly on a step boundary.
    A non-integer ratio is rejected rather than rounded, since rounding it would silently change
    the sample rate from the configured one.
    \item \textbf{The hold is constant across every stage of a substep.} Inside a tick the
    right-hand side is then again a smooth function of $(t,\vec{y})$, so the integrator's order is
    preserved; every discontinuity introduced by \autoref{eq:zoh} lands on a step boundary and
    never inside one. This is what makes the mode a true zero-order hold and not an approximation
    of one.
    \item \textbf{The sample is taken at the tick's start}, $(t_k,\vec{y}(t_k))$ --- the causally
    correct order for a flight computer: measure, compute, apply until the next tick.
\end{enumerate}

Two properties of a sampled run should be stated plainly, because neither is a defect awaiting
correction. First, such a run is a \emph{different dynamical system} from the continuously
evaluated one, not a more accurate solution of the same one: the two are not to be compared as
though one refined the other, and the comparison that does mean something is ``this actuator at
this sample rate'' against ``an ideal continuous actuator''. Second, it is substantially more
expensive --- small fixed steps for the whole flight, against the adaptive scheme's far smaller
step count --- which is why it is opt-in rather than the default. $f_s=0$ is that default, and with
it the mode is structurally absent rather than merely inert.

This mode supplies the place to run a controller with memory, and nothing else: no such
controller, and no state-coupled turbulence model, is written here.

\end{aitext}

\newpage
\begin{aitext}
\section{Backlog and future work}\label{sec:backlog}
The complement of \autoref{sec:assumptions}: that section records what the present model
simplifies, this one what is deliberately left undone. Each entry says what the item is and ---
more usefully --- what has to exist before it can be written down properly, so that a reader can
tell a missing derivation from a missing input.

\begin{itemize}
    \item \textbf{Commanded gimbal control (direction/guidance).} $\delta$ is, for now, deliberately
    kept as a constant input (\autoref{sec:thrust}) --- the point of the current pass was to see how
    a fixed manufacturing-type misalignment behaves (\autoref{sec:analysis_misaligned}) before
    adding any control authority. A \emph{commanded}, time-varying $\delta(t)$ driven by a control
    law is the next step once that is wanted, and is explicitly a 3D-model item: steering needs an
    out-of-plane degree of freedom this 2D closure does not have. \emph{Waiting on:} a stated
    guidance objective (pitch programme? gravity turn? impact point?), actuator rate and deflection
    limits, a sensor model to close the loop on, and the 3D model itself. A \emph{different} control
    loop --- drag modulation via a feedback-controlled air brake, targeting a commanded apogee
    during unpowered coast --- is already implemented and stays 2D-only (\autoref{sec:airbrake}):
    it needs no out-of-plane freedom, only an additional switchable drag term. Its control law is a
    continuous, memoryless proportional one, not a sampled-data controller with integral/derivative
    memory --- the different integration strategy such a controller needs now exists
    (\autoref{sec:sampled}), but the law itself is still future work (see the sampled-data item
    below). Two smaller gaps in the same loop belong here rather than there: the brake has no
    retraction condition and no deployment-rate limit (\autoref{sec:airbrake_trigger}), and it
    perturbs $A_x$ alone, leaving $A_z$, $A_m$ and $l_1$ untouched --- which the wing-body item
    below is what would fix properly.
    \item \textbf{Wing-body aerodynamics: analytic CoP, airframe, air brake.} Replace the
    tabulated, linearized or constant coefficients of \autoref{sec:aero} with a wing-body
    (Barrowman-style) build-up. Conceptually: the nose, the body, and each fin each contribute
    their own normal-force slope $C_{N\alpha,i}$ and their own centre-of-pressure location
    $X_i$ (a slender
    nose/body contributes a small, often near-zero net slope away from the nose in the classical
    linear treatment, aside from any boat-tail correction; each fin's contribution uses a
    corrected 2D lift-curve slope, an interference factor between fin and body, and an
    effective aspect ratio from its planform), and the assembly's centre of pressure is the
    $C_{N\alpha}$-weighted average of the component locations,
    $X_{cp}=\sum_iC_{N\alpha,i}X_i/\sum_iC_{N\alpha,i}$ --- so that $l_1(\alpha,M)$ is
    \emph{computed} from the airframe rather than read from a table, and an air brake enters as a
    switchable change to $C_A$ and to the CoP. That last point of entry is already fixed by the
    structure of \autoref{sec:aero}: whatever a control surface does, it does by perturbing
    $C_x,C_z,C_m$ or $l_1$, so what this item is waiting on is the airframe and the deflection law
    --- not somewhere to put them. The air brake of \autoref{sec:airbrake} is that entry point
    exercised \emph{without} an airframe behind it: a bare drag area added to $A_x$, which is
    precisely why it needs none of the build-up above and why it cannot say anything about the
    normal force or the centre of pressure. The original Barrowman formulas are subsonic
    linear-theory results and would need a Mach correction (Prandtl--Glauert or, in the linear
    supersonic regime, Busemann theory) to extend past low subsonic speed.
    \emph{Waiting on:} an actual airframe --- nose shape, body diameter/length and any boat-tail,
    and fin planform (root/tip chord, span, sweep, thickness, count) --- which this document does
    not yet fix anywhere.
    \item \textbf{Dynamic stability, quantitatively.} The framework is now in the document
    (\autoref{sec:analysis_dynamic}): a pitch-damping derivative $c_{m,q}$ linearises into a
    damped short-period oscillator with natural frequency $\omega_n$ and damping ratio $\zeta$.
    What is missing is a value or model for $c_{m,q}$ itself. \emph{Waiting on:} either a measured
    value (forced-oscillation wind-tunnel testing, or unsteady CFD) or a full distributed-force
    derivation from the airframe geometry of the previous item.
    \item \textbf{Propellant grain geometry behind $I_{m,g}(t)$ and $l_{g,m}(t)$.} The physical
    mechanism and the standard grain families (cylindrical/BATES, tubular, end-burning,
    star/finocyl) are now elaborated in \autoref{sec:mass}; what remains is choosing one and fixing
    its actual dimensions (bore/outer diameter, length, and, for star grains, the point geometry),
    from which $A_b(w)$, $l_{g,m}(t)$ and $I_{m,g}(t)$ follow by direct integration. Picking this
    up is the next conversation, not something to guess at here.
    \item \textbf{Component-based estimate of the dry rocket's mass properties.} Today the dry
    rocket's mass $m_r$, centre of gravity $l_{g,r}$ and moment of inertia $I_{r,g}$
    (\autoref{sec:mass}) are entered by hand. An OpenRocket-style estimator would build them from
    components --- nose cone, body tubes, fins, internal parts --- each with its geometry and
    material density, combining masses, centroids and own moments of inertia by the same
    mass-weighted-centroid and parallel-axis rules as \autoref{eq:mass_cog} and
    \autoref{eq:mass_inertia_parts}. \emph{Waiting on:} the airframe geometry of the wing-body item
    above and a materials table.
    \item \textbf{Vertical wind.} Two candidate families exist (terrain-slope/mass-continuity
    orographic lift, or a convective/thermal parameterisation, \autoref{sec:wind} > Vertical
    wind), each needing a data source (terrain elevation, or surface heat-flux/insolation data)
    this project does not currently have. Deliberately on hold pending research into which family,
    if either, is worth the added dependency --- not to be implemented before that research is
    done.
    \item \textbf{CFD/data-extraction pipeline for $C_A,C_N,C_m$.} \autoref{sec:aero} now states
    plainly that these coefficients, and $l_1$, are supplied rather than derived in this document,
    and that the highest-fidelity of its three ways of supplying them --- the table model --- reads
    them from a database (wind-tunnel testing, CFD sweeps, or a semi-empirical panel-method tool).
    What remains open is the pipeline that actually populates that database for a specific
    airframe --- the CFD run matrix or wind-tunnel campaign itself --- which is outside the scope
    of this document.
    \item Out-of-plane motion: crosswind weathercocking (\autoref{sec:wind}) and the lateral
    Coriolis deflection (\autoref{sec:noninertial}). Both are discarded by the same 2D closure and
    both are lifted by the same future work, a full 6DOF model --- not by anything that can be
    done inside this document.
    \item A sampled-data control \emph{law} --- one with integral/derivative memory, running on the
    fixed-rate, zero-order-hold loop --- for the air brake of \autoref{sec:airbrake}, or for any
    future controller. The integration strategy this needs is \emph{no longer} missing
    (\autoref{sec:sampled}, which also states the two properties of a sampled run that are not
    defects to be corrected: it is a different dynamical system rather than a refinement of the
    continuous one, and it is substantially more expensive). This item has therefore gone from
    blocked to merely unwritten. \emph{Waiting on:} the law itself --- which signal it acts on, its
    gains, and the saturation/anti-windup treatment of its integral term --- none of which this
    document fixes. The present proportional law (\autoref{sec:airbrake_law}) remains re-evaluated
    fresh at every call, and needs none of this.
    \item \textbf{State-coupled turbulence.} The Dryden spectrum of \autoref{sec:wind} >
    Turbulence is driven by a fixed reference airspeed, not by the rocket's own instantaneous one
    --- the gold-standard formulation, a filter driven by the true airspeed and carrying its own
    internal state, is the one a flight-test comparison would want. This has the same root cause as
    the item above and is unblocked by the same mechanism: a filter state is exactly the recursive
    per-tick quantity \autoref{eq:zoh} carries. \emph{Waiting on:} the filter's own discretisation,
    and a decision on whether drawing noise per tick turns the flight into a stochastic
    differential equation whose results then depend on the sample rate itself --- a genuine
    modelling question, not an implementation detail.
    \item A refined launch-rail model with friction and a tip-off transient as the last button
    leaves the rail (\autoref{sec:rail}), and a recovery model with the canopy's own opening
    transient, inflation shock and pendulum swing under the risers, rather than the current
    instantaneous-deployment, frozen-attitude treatment (\autoref{sec:recovery}).
\end{itemize}
\end{aitext}

\end{document}
